%%\documentclass[a4paper,14pt,twoside,justified, symmetric]{tufte-book} % Use the tufte-book class which in turn uses the tufte-common class
%\documentclass{book}
%\usepackage[margin=0.7in]{geometry}
%
%%%%%%%%%%%%%%%%%%%PAGE SIZE%%%%%%
%%\geometry{textwidth=.58\paperwidth} % size of text column
%%\geometry{
%%  left=21.8mm, % left margin
%%   textwidth=115mm
%%%   , % main text block
%%%  marginparsep=8.2mm, % gutter between main text block and margin notes
%%%  marginparwidth=49.4mm % width of margin notes
%%}
%
% %%%%%%%%%%%%%%%%%%PAGE SIZE%%%%%%
%\usepackage{graphicx}
%\usepackage{mathpazo}
%\usepackage{lscape}
%\usepackage[usenames,dvipsnames]{xcolor}
%
%\usepackage{dingbat}
%\usepackage{wrapfig}
%\usepackage{dirtytalk}
%\usepackage{pifont}% http://ctan.org/pkg/pifont
%\newcommand{\cmark}{\ding{51}}%
%\newcommand{\xmark}{\ding{55}}%
%%\usepackage{fourier}
%%\usepackage{marvosym}\usepackage{pifont}
%%\usepackage{url}
%\usepackage{tabu}
%\usepackage{modiagram}
%\usepackage{amsmath}
%\usepackage{setspace}
%\usepackage{afterpage}
%\usepackage[version=4]{mhchem}
%\usepackage{fontawesome}
%\usepackage{soul}
%\sethlcolor{green}
%\usepackage{amsmath}
%\usepackage{fonttable}
%\usepackage{paracol}
%\usepackage{chemfig}
%\usetikzlibrary{3d}
%\usepackage[ugly]{nicefrac}
%\usepackage{cutwin}
%%\usepackage{pstricks}
%
%\usepackage{pgfplots}\pgfplotsset{compat=newest}
%
%\usepackage{movie15}
%\usepackage{hyperref}
%
%
%\graphicspath{ {images/} }
%\hypersetup{colorlinks} % Comment this line if you don't wish to have colored links
%\usepackage{microtype} % Improves character and word spacing
%\usepackage{lipsum} % Inserts dummy text
%\usepackage{booktabs} % Better horizontal rules in tables
%%%%%%%COLORS
%\definecolor{olive}{rgb}{0.3, 0.4, .1}
%\definecolor{fore}{RGB}{249,242,215}
%\definecolor{back}{RGB}{51,51,51}
%\definecolor{title}{RGB}{255,0,90}
%\definecolor{dgreen}{rgb}{0.,0.6,0.}
%\definecolor{gold}{rgb}{1.,0.84,0.}
%\definecolor{JungleGreen}{cmyk}{0.99,0,0.52,0}
%\definecolor{BlueGreen}{cmyk}{0.85,0,0.33,0}
%\definecolor{RawSienna}{cmyk}{0,0.72,1,0.45}
%\definecolor{Magenta}{cmyk}{0,1,0,0}
%%%%%%%COLORS
%\usepackage{graphicx} % Needed to insert images into the document
%\graphicspath{{graphics/}} % Sets the default location of pictures
%\setkeys{Gin}{width=\linewidth,totalheight=\textheight,keepaspectratio} % Improves figure scaling
%\usepackage{fancyvrb} % Allows customization of verbatim environments
%\fvset{fontsize=\normalsize} % The font size of all verbatim text can be changed here
%\newcommand{\hangp}[1]{\makebox[0pt][r]{(}#1\makebox[0pt][l]{)}} % New command to create parentheses around text in tables which take up no horizontal space - this improves column spacing
%\newcommand{\hangstar}{\makebox[0pt][l]{*}} % New command to create asterisks in tables which take up no horizontal space - this improves column spacing
%\usepackage{xspace} % Used for printing a trailing space better than using a tilde (~) using the \xspace command
%\newcommand{\monthyear}{\ifcase\month\or January\or February\or March\or April\or May\or June\or July\or August\or September\or October\or November\or December\fi\space\number\year} % A command to print the current month and year
%\newcommand{\openepigraph}[2]{ % This block sets up a command for printing an epigraph with 2 arguments - the quote and the author
%
%\sffamily\large
%\begin{doublespace}
%\noindent\allcaps{#1}\\ % The quote
%\noindent\allcaps{#2} % The author
%\end{doublespace}
% }
%\newcommand{\blankpage}{\newpage\hbox{}\thispagestyle{empty}\newpage} % Command to insert a blank page
%\usepackage{units} % Used for printing standard units
%\newcommand{\hlred}[1]{\textcolor{Maroon}{#1}} % Print text in maroon
%%\newcommand{\hangleft}[1]{\makebox[0pt][r]{#1}} % Used for printing commands in the index, moves the slash left so the command name aligns with the rest of the text in the index 
%\newcommand{\hairsp}{\hspace{1pt}} % Command to print a very short space
%\newcommand{\ie}{\textit{i.\hairsp{}e.}\xspace} % Command to print i.e.
%\newcommand{\eg}{\textit{e.\hairsp{}g.}\xspace} % Command to print e.g.
%\newcommand{\na}{\quad--} % Used in tables for N/A cells
%\newcommand{\measure}[3]{#1/#2$\times$\unit[#3]{pc}} % Typesets the font size, leading, and measure in the form of: 10/12x26 pc.
%\newcommand{\tuftebs}{\symbol{'134}} % Command to print a backslash in tt type in OT1/T1
%\providecommand{\XeLaTeX}{X\lower.5ex\hbox{\kern-0.15em\reflectbox{E}}\kern-0.1em\LaTeX}
%\newcommand{\tXeLaTeX}{\XeLaTeX\index{XeLaTeX@\protect\XeLaTeX}} % Command to print the XeLaTeX logo while simultaneously adding the position to the index
%\newcommand{\doccmdnoindex}[2][]{\texttt{\tuftebs#2}} % Command to print a command in texttt with a backslash of tt type without inserting the command into the index
%\newcommand{\doccmddef}[2][]{\hlred{\texttt{\tuftebs#2}}\label{cmd:#2}\ifthenelse{\isempty{#1}} % Command to define a command in red and add it to the index
%{ % If no package is specified, add the command to the index
%\index{#2 command@\protect\hangleft{\texttt{\tuftebs}}\texttt{#2}}% Command name
%}
%{ % If a package is also specified as a second argument, add the command and package to the index
%\index{#2 command@\protect\hangleft{\texttt{\tuftebs}}\texttt{#2} (\texttt{#1} package)}% Command name
%\index{#1 package@\texttt{#1} package}\index{packages!#1@\texttt{#1}}% Package name
%}}
%\newcommand{\doccmd}[2][]{% Command to define a command and add it to the index
%\texttt{\tuftebs#2}%
%\ifthenelse{\isempty{#1}}% If no package is specified, add the command to the index
%{%
%\index{#2 command@\protect\hangleft{\texttt{\tuftebs}}\texttt{#2}}% Command name
%}
%{%
%\index{#2 command@\protect\hangleft{\texttt{\tuftebs}}\texttt{#2} (\texttt{#1} package)}% Command name
%\index{#1 package@\texttt{#1} package}\index{packages!#1@\texttt{#1}}% Package name
%}}
%% A bunch of new commands to print commands, arguments, environments, classes, etc within the text using the correct formatting
%\newcommand{\docopt}[1]{\ensuremath{\langle}\textrm{\textit{#1}}\ensuremath{\rangle}}
%\newcommand{\docarg}[1]{\textrm{\textit{#1}}}
%\newenvironment{docspec}{\begin{quotation}\ttfamily\parskip0pt\parindent0pt\ignorespaces}{\end{quotation}}
%\newcommand{\docenv}[1]{\texttt{#1}\index{#1 environment@\texttt{#1} environment}\index{environments!#1@\texttt{#1}}}
%\newcommand{\docenvdef}[1]{\hlred{\texttt{#1}}\label{env:#1}\index{#1 environment@\texttt{#1} environment}\index{environments!#1@\texttt{#1}}}
%\newcommand{\docpkg}[1]{\texttt{#1}\index{#1 package@\texttt{#1} package}\index{packages!#1@\texttt{#1}}}
%\newcommand{\doccls}[1]{\texttt{#1}}
%\newcommand{\docclsopt}[1]{\texttt{#1}\index{#1 class option@\texttt{#1} class option}\index{class options!#1@\texttt{#1}}}
%\newcommand{\docclsoptdef}[1]{\hlred{\texttt{#1}}\label{clsopt:#1}\index{#1 class option@\texttt{#1} class option}\index{class options!#1@\texttt{#1}}}
%\newcommand{\docmsg}[2]{\bigskip\begin{fullwidth}\noindent\ttfamily#1\end{fullwidth}\medskip\par\noindent#2}
%\newcommand{\docfilehook}[2]{\texttt{#1}\index{file hooks!#2}\index{#1@\texttt{#1}}}
%\newcommand{\doccounter}[1]{\texttt{#1}\index{#1 counter@\texttt{#1} counter}}
%\usepackage{makeidx} % Used to generate the index
%\makeindex % Generate the index which is printed at the end of the document
%% This block contains a number of shortcuts used throughout the book
%\newcommand{\vdqi}{\textit{VDQI}\xspace}
%\newcommand{\ei}{\textit{EI}\xspace}
%\newcommand{\ve}{\textit{VE}\xspace}
%\newcommand{\be}{\textit{BE}\xspace}
%\newcommand{\VDQI}{\textit{The Visual Display of Quantitative Information}\xspace}
%\newcommand{\EI}{\textit{Envisioning Information}\xspace}
%\newcommand{\VE}{\textit{Visual Explanations}\xspace}
%\newcommand{\BE}{\textit{Beautiful Evidence}\xspace}
%\newcommand{\TL}{Tufte-\LaTeX\xspace}
%
%
%%%----------------------------------------------------------------------------------------
%%%	BOOK META-INFORMATION
%%%----------------------------------------------------------------------------------------
%%\title{Chemistry For Allied Health: an OpenTextbook} % Title of the book
%%\author[Prof. Daniel Torres, \today]{Daniel Torres} % Author
%%\publisher{BMCC-CUNY} % Publisher
%%%----------------------------------------------------------------------------------------
%
%
%%%%%%%%%%%%%%%%%%%%%%%%%Chapter Numbers
% \def\chpcolor{blue!30!black}
%\def\chpcolortxt{black!90}
%\def\sectionfont{\sffamily\LARGE}
%\setcounter{secnumdepth}{2}
%\makeatletter
%%Section:
%\def\@sectionstrut{\vrule\@width\z@\@height12.5\p@}
%\def\@makesectionhead#1{%
%  {\par\vspace{20pt}%
%   \parindent 0pt\raggedleft\sectionfont
%   \colorbox{\chpcolor}{%
%     \parbox[t]{90pt}{\color{white}\@sectionstrut\@depth4.5\p@\hfill
%       \ifnum\c@secnumdepth>\z@\thesection\fi}%
%   }%
%   \begin{minipage}[t]{\dimexpr\textwidth-90pt-2\fboxsep\relax}
%   \color{\chpcolortxt}\@sectionstrut\hspace{5pt}#1
%   \end{minipage}\par
%   \vspace{10pt}%
%  }
%}
%\def\section{\@afterindentfalse\secdef\@section\@ssection}
%\def\@section[#1]#2{%
%  \ifnum\c@secnumdepth>\m@ne
%    \refstepcounter{section}%
%    \addcontentsline{toc}{section}{\protect\numberline{\thesection}#1}%
%  \else
%    \phantomsection
%    \addcontentsline{toc}{section}{#1}%
%  \fi
%  \sectionmark{#1}%
%  \if@twocolumn
%    \@topnewpage[\@makesectionhead{#2}]%
%  \else
%    \@makesectionhead{#2}\@afterheading
%  \fi
%}
%\def\@ssection#1{%
%  \if@twocolumn
%    \@topnewpage[\@makesectionhead{#1}]%
%  \else
%    \@makesectionhead{#1}\@afterheading
%  \fi
%}
%\makeatother
%%%%%%%%%%%%%%%%%%%%%%%%%Chapter Numbers
%
%%%%%%%BOXES SAMPLE PROBLEM     
%\usepackage[many]{tcolorbox}
%\definecolor{myblue}{RGB}{0,163,243}
%\tcbset{mystyle/.style={
%  breakable,
%  enhanced,
%  outer arc=0pt,
%  arc=0pt,
%  colframe=black!45,
%  colback=black!45,
%  attach boxed title to top left,
%  boxed title style={
%    colback=black!40,
%    outer arc=0pt,
%    arc=0pt,
%    },
%  title=Example~\thetcbcounter,
%  fonttitle=\sffamily
%  }
%}
%\tcbset{mystyle2/.style={
%  breakable,
%  enhanced,
%  outer arc=0pt,
%  arc=0pt,
%  colframe=black!45,
%  colback=white,
%  attach boxed title to top left,
%  boxed title style={
%    colback=white,
%    outer arc=0pt,
%    arc=0pt,
%    },
%  fonttitle=\sffamily
%  }
%}
%\newtcolorbox[auto counter]{example}[1][]{
%  mystyle,
%  colback=white,
%  rightrule=0pt,
%  toprule=0pt,
%  title=Sample Problem \thetcbcounter,
%  }
%  \newtcolorbox{work}[1][]{
%  mystyle2,
%  colback=white,
%  rightrule=0pt,
%  toprule=0pt,  
%  bottomrule=0pt,
%  rightrule=0pt,
%  }
%  \newtcolorbox[auto counter]{problem}[1][]{
%  mystyle,
%  colback=white,
%  rightrule=0pt,
%  toprule=0pt,
%  title=Chapter Problems,
%  }
%%%%%%%BOXES SAMPLE PROBLEM
%\newcommand\textlcsc[1]{\textsc{\MakeLowercase{#1}}}   %FORCE CAPS
%%%%%%%COLOED CIRCLES
%\newcommand*\circled[1]{\tikz[baseline=(char.base)]{%
%            \node[shape=circle,fill=blue!30!black,,draw,inner sep=2pt] (char) {#1};}}
%\usepackage{enumitem}
%%%%%%%COLOED CIRCLES
%\usepackage{textcomp}
%\usepackage{lettrine}
%\usepackage{lipsum}
%\usepackage{lmodern}
%\usepackage{tcolorbox}
%\usepackage[framemethod=TikZ]{mdframed}
%%%%PROBLEMS BOX
%\mdfdefinestyle{problems}{%
%    linecolor=black,
%    outerlinewidth=2pt,
%    roundcorner=20pt,
%    innertopmargin=\baselineskip,
%    innerbottommargin=\baselineskip,
%    innerrightmargin=20pt,
%    innerleftmargin=20pt,
%    backgroundcolor=gray!10!white}
%%%%PROBLEMS BOX 
%%%%PROBLEMS BOX 2
%\usepackage{blindtext}
%\usepackage{amsmath,amssymb}
%\usepackage{lmodern}
%\usepackage{tcolorbox}
%\makeatletter
%\newtcolorbox{mytcbox}[2][]{%
%enhanced, 
%breakable,
%colback=white,
%colframe=blue!30!black,
%attach boxed title to top right={yshift=-2pt}, title={#2},
%boxed title size=standard,
%boxrule=0pt,
%boxed title style={%
%    sharp corners, 
%    rounded corners=northeast, 
%    colback=tcbcol@frame, 
%    boxrule=0pt},
%sharp corners=north,
%overlay unbroken={%
%    \path[fill=tcbcol@back] 
%        ([xshift=2pt]title.south west) 
%        to[out=180, in=0] ([xshift=-3.5cm]title.west)--
%        (title.west-|frame.west) |- 
%        ([xshift=2pt]title.south west)--cycle;
%    \path[fill=tcbcol@frame] (title.south west) 
%        to[out=180, in=0] ([xshift=-1.5cm]title.west)--
%        (title.west-|frame.west)
%        [rounded corners=\kvtcb@arc] |- 
%        (title.north-|frame.north) 
%        [sharp corners] -| (title.south west);
%    \draw[line width=.5mm, rounded corners=\kvtcb@arc, 
%        tcbcol@frame] 
%        (title.north east) rectangle 
%        (frame.south west);
%}, 
%overlay first={%
%    \path[fill=tcbcol@back] 
%        ([xshift=2pt]title.south west) 
%        to[out=180, in=0] ([xshift=-1.5cm]title.west)--
%        (title.west-|frame.west) |- 
%        ([xshift=2pt]title.south west)--cycle;
%    \path[fill=tcbcol@frame] (title.south west) 
%        to[out=180, in=0] ([xshift=-1.5cm]title.west)--
%        (title.west-|frame.west)
%        [rounded corners=\kvtcb@arc] |- 
%        (title.north-|frame.north) 
%        [sharp corners] -| (title.south west);
%    \draw[line width=.5mm, rounded corners=\kvtcb@arc, 
%        tcbcol@frame] 
%        (frame.south west) |- (title.north) -| 
%        (frame.south east);
%}, 
%overlay middle={%
%    \draw[line width=.5mm, tcbcol@frame] 
%    (frame.north west)--(frame.south west) 
%    (frame.north east)--(frame.south east);
%}, 
%overlay last={%
%    \draw[line width=.5mm, rounded corners=\kvtcb@arc, 
%        tcbcol@frame] 
%        (frame.north west) |- (frame.south) -|
%        (frame.north east);
%}, 
%#1
%}
%\makeatother
%%%%PROBLEMS BOX 2
%%%%DISCUSS BOX
%\makeatletter
%
%\def\tcb@shadow@lifted#1#2#3#4{%
%  \path[fill,rounded corners=\tcb@outer@arc,#4]
%       ([xshift=#1+#3,yshift=#2+#3]frame.south west)
%    .. controls ([yshift=\dimexpr#3]frame.south) ..
%       ([xshift=-#1-#3,yshift=#2+#3]frame.south east)
%    -- ([xshift=-#1-#3,yshift=#2-#3]frame.north east)
%    -- ([xshift=#1+#3,yshift=#2-#3]frame.north west)
%    -- cycle;
%}
%
%\tcbset{
%  lifted shadow/.style args={#1#2#3#4}{shad@w app={%
%    \begin{scope}[#4]%
%      \tcb@shadow@lifted{#1}{#2}{\dimexpr-4\dimexpr#3}{opacity=0.01}%
%      \tcb@shadow@lifted{#1}{#2}{\dimexpr-3\dimexpr#3}{opacity=0.02}%
%      \tcb@shadow@lifted{#1}{#2}{\dimexpr-2\dimexpr#3}{opacity=0.04}%
%      \tcb@shadow@lifted{#1}{#2}{\dimexpr-#3}{opacity=0.07}%
%      \tcb@shadow@lifted{#1}{#2}{0pt}{opacity=0.11}%
%      \tcb@shadow@lifted{#1}{#2}{\dimexpr+#3}{opacity=0.11}%
%      \tcb@shadow@lifted{#1}{#2}{\dimexpr+2\dimexpr#3}{opacity=0.07}%
%      \tcb@shadow@lifted{#1}{#2}{\dimexpr+3\dimexpr#3}{opacity=0.04}%
%      \tcb@shadow@lifted{#1}{#2}{\dimexpr+4\dimexpr#3}{opacity=0.02}%
%      \tcb@shadow@lifted{#1}{#2}{\dimexpr+5\dimexpr#3}{opacity=0.01}%
%  \end{scope}}},%
%  drop lifted shadow/.style={lifted shadow={1.5mm}{-1.5mm}{0.12mm}{#1}},
%  drop lifted shadow/.default={black!50!white},%
%  drop heavy lifted shadow/.style={lifted shadow={2mm}{-3mm}{0.16mm}{#1}},
%  drop heavy lifted shadow/.default={black!50!white},%
%}
%
%\makeatother
%%%%DISCUSS BOX
%
%%BOX FOR EQUATIONS
%\usetikzlibrary{calc}
%\newcommand*{\boxcolor}{blue!30!black}
%\makeatletter
%\renewcommand{\boxed}[1]{\textcolor{\boxcolor}{%
%\tikz[baseline={([yshift=-1ex]current bounding box.center)}] \node [rectangle, minimum width=1ex,rounded corners,draw] {\normalcolor\m@th$\displaystyle#1$};}}
% \makeatother
%
%%BOX FOR EQUATIONS
%
%
%\usepackage{amsmath}
%\usepackage[makeroom]{cancel}
%
%
%\setlength{\columnsep}{.5cm}
%
%
%
%\usepackage{exsheets}
%\SetupExSheets{counter-format=se.qu}
%
%\NewTasks[style=multiplechoice]{choices}[\choice]
%\newcommand*\correct{\checkedchoicebox}
%
%%TABLES
%\usepackage{xcolor,colortbl}
%\newcommand{\mc}[2]{\multicolumn{#1}{c}{#2}}
%\usepackage{import}
%\usepackage{subfiles}
%%YELLOW NOTES
%\setlength{\parskip}{0ex}
%\setlength{\parindent}{0ex}
%
%\newlength{\yellownotewidth}
%\setlength{\yellownotewidth}{2.5cm}
%\newlength{\yellownoteheight}
%\setlength{\yellownoteheight}{2.5cm}
%%   -   -   -   -   -   -   -   -   -   -   -   -
%% Yellow note...
%%   -   -   -   -   -   -   -   -   -   -   -   -
%\newcommand{\resizeableyellownote}[3]{
%\setlength{\yellownotewidth}{#1cm}
%\setlength{\yellownoteheight}{#2cm}
%\marginpar{
%    \vspace{-0.5\yellownoteheight}
%        \begin{center}
%        \begin{tikzpicture}
%            \draw[white,fill=gray!25,opacity=0.75,shift={(-0.125,-0.125)}] 
%                (0,0) rectangle (\yellownotewidth,\yellownoteheight);
%            \draw[fill=yellow!35] (0,0) rectangle (\yellownotewidth,\yellownoteheight);
%            \draw[opacity=0.45,fill=gray!50] (0.7\yellownotewidth,0) -- 
%                (0.9\yellownotewidth,0.45) -- (\yellownotewidth,0.4) -- cycle;
%            \node[blue,below] at (0.5\yellownotewidth,\yellownoteheight) {
%                \begin{minipage}{\yellownotewidth-1em}
%                    \scriptsize\sf#3
%                \end{minipage}
%            };
%        \end{tikzpicture}
%        \end{center}
%        \vspace{0.5\yellownoteheight}
%    }
%}
%%YELLOW NOTES
%%%%% CALCULATOR KEY
%\usepackage[T1]{fontenc}
%\usepackage[utf8]{inputenc}
%\usepackage{tikz}
%\usetikzlibrary{shadows}
%\newcommand*\keystroke[1]{%
%  \tikz[baseline=(key.base)]
%    \node[%
%      draw,
%      fill=white,
%      drop shadow={shadow xshift=0.25ex,shadow yshift=-0.25ex,fill=black,opacity=0.75},
%      rectangle,
%      rounded corners=2pt,
%      inner sep=1pt,
%      line width=0.5pt,
%      font=\scriptsize\sffamily
%    ](key) {#1\strut}
%  ;
%}
%%%%% CALCULATOR KEY
%
%%HIGHLIGHT IN MATH
%\newcommand{\hlmath}[1]{%
%  \colorbox{black!20}{$\displaystyle#1$}}
%  %HIGHLIGHT IN MATH
%%EMAIL SEND
%\usepackage{etoolbox}
%\makeatletter
%\newcommand\myemail[3]{%                %\newcommand\tpj@compose@mailto[3]{%
%\edef\@tempa{mailto:#1?subject=#2 }%
%\edef\@tempb{\expandafter\html@spaces\@tempa\@empty}%
%\href{\@tempb}{#3}}
%\catcode\%=11
%\def\html@spaces#1 #2{#1%20\ifx#2\@empty\else\expandafter\html@spaces\fi#2}
%\catcode\%=14
%\makeatother
%\usepackage{hyperref}
%%EMAIL SEND
%%DECORATION LINE
%\usepackage[object=vectorian]{pgfornament} %%  http://altermundus.com/pages/tkz/ornament/index.html
%\newcommand{\sectionlinetwo}[2]{%
%  \nointerlineskip \vspace{0.5\baselineskip}\hspace{\fill}
%  {\color{#1}
%    \resizebox{0.8\linewidth}{2ex}
%    {{%
%    {\begin{tikzpicture}
%    \node  (C) at (0,0) {};
%    \node (D) at (9,0) {};
%    \path (C) to [ornament=#2] (D);
%    \end{tikzpicture}}}}}%
%    \hspace{\fill}
%    \par\nointerlineskip \vspace{.5\baselineskip}
%  }
%%DECORATION LINE
%% Periodic tabe
%\usetikzlibrary{shapes,calc}
%
%\ifpdf
%  %
%\else
%  % Implement Outline text using pstricks if regular LaTeX->dvi->ps->pdf route
%  \usepackage{pst-all}
%\fi
%% Periodic tabe
%% QUOTES
%\usepackage{ifxetex,ifluatex}
%\usepackage{etoolbox}
%\usepackage{xcolor}
%\usepackage{framed}
%\newif\ifxetexorluatex
%\ifxetex
%  \xetexorluatextrue
%\else
%  \ifluatex
%    \xetexorluatextrue
%  \else
%    \xetexorluatexfalse
%  \fi
%\fi
%%
%\ifxetexorluatex%
%  \usepackage{fontspec}
%  \usepackage{libertine} % or use \setmainfont to choose any font on your system
%  \newfontfamily\quotefont[Ligatures=TeX]{Linux Libertine O} % selects Libertine as the quote font
%\else
%  \usepackage[utf8]{inputenc}
%  \usepackage[T1]{fontenc}
%  \usepackage{libertine} % or any other font package
%  \newcommand*\quotefont{\fontfamily{LinuxLibertineT-LF}} % selects Libertine as the quote font
%\fi
%\newcommand*\quotesize{60} % if quote size changes, need a way to make shifts relative
%% Make commands for the quotes
%\newcommand*{\openquote}
%   {\tikz[remember picture,overlay,xshift=-4ex,yshift=-2.5ex]
%   \node (OQ) {\quotefont\fontsize{\quotesize}{\quotesize}\selectfont``};\kern0pt}
%
%\newcommand*{\closequote}[1]
%  {\tikz[remember picture,overlay,xshift=4ex,yshift={#1}]
%   \node (CQ) {\quotefont\fontsize{\quotesize}{\quotesize}\selectfont''};}
%% select a colour for the shading
%%\colorlet{shadecolor}{Azure}
%\newcommand*\shadedauthorformat{\emph} % define format for the author argument
%
%% Now a command to allow left, right and centre alignment of the author
%\newcommand*\authoralign[1]{%
%  \if#1l
%    \def\authorfill{}\def\quotefill{\hfill}
%  \else
%    \if#1r
%      \def\authorfill{\hfill}\def\quotefill{}
%    \else
%      \if#1c
%        \gdef\authorfill{\hfill}\def\quotefill{\hfill}
%      \else\typeout{Invalid option}
%      \fi
%    \fi
%  \fi}
%% wrap everything in its own environment which takes one argument (author) and one optional argument
%% specifying the alignment [l, r or c]
%%
%\newenvironment{shadequote}[2][l]%
%{\authoralign{#1}
%\ifblank{#2}
%   {\def\shadequoteauthor{}\def\yshift{-2ex}\def\quotefill{\hfill}}
%   {\def\shadequoteauthor{\par\authorfill\shadedauthorformat{#2}}\def\yshift{2ex}}
%\begin{snugshade}\begin{quote}\openquote}
%{\shadequoteauthor\quotefill\closequote{\yshift}\end{quote}\end{snugshade}}
%% QUOTES
%% FRAMED BOX
%\usepackage{mwe}
%\usetikzlibrary{shadows}
%\definecolor{mygreen}{rgb}{0.61,0.73,0.35}
%\mdfdefinestyle{mystyle}{%
%    innerlinewidth = 0.5pt,
%    outerlinewidth = 2pt,
%    linecolor = mygreen,
%    tikzsetting = {draw = white, line width = 0.5pt},
%    roundcorner=15pt,
%    linecolor=mygreen,
%    linewidth=2pt,
%    topline=true,
%    frametitleaboveskip=1.5\dimexpr-\ht\strutbox\relax,
%}
%
%\newenvironment{definition}[1][]{%
%\ifstrempty{#1}%
%    {\mdfsetup{%
%    style = mystyle,
%    }}%
%    {\mdfsetup{%
%    style = mystyle,
%    frametitle={%
%    \tikz{[baseline=(current bounding box.east),outer sep=0pt]
%    \node[draw = white, line width = 2pt, text = white, anchor=east,rectangle,
%    fill=mygreen, rounded corners, drop shadow]
%    {\strut #1};
%    }}}}%
%    \begin{mdframed}[]\relax%
%}{\end{mdframed}}
%%FRAMED BOX
%
%%% NEW PROBLEMS BOX
%%\usepackage{tikz}
%%\usetikzlibrary{calc}
%%
%%\xdefinecolor{mycolor3}{RGB}{62,96,111} % Neutral Blue
%%\colorlet{bancolor}{mycolor3}
%%
%%\def\bancolor{mycolor3}
%%\newenvironment{mybox3}[3][]{%
%%    \begin{tikzpicture}[#1]%
%%        \def\myboxname{#3}%
%%        \node [draw,inner sep=1.5ex,text width=#2]% good options: minimum height, minimum width
%%            (BOXCONTENT) \bgroup\rule{0pt}{3ex}\ignorespaces
%%}{%
%%        \egroup;
%%        \node [right,inner xsep=1em,fill=bancolor!75,outer sep=0pt,text height=2ex,text depth=.5ex] (BOXNAME) 
%%            at ([shift={(-1em,0pt)}]BOXCONTENT.north west) {\myboxname};
%%        \fill[bancolor] (BOXNAME.north east) -- +(-1em,1em) -- +(-1em,0) -- cycle;
%%        \fill[bancolor] (BOXNAME.south west) -- +(1em,-1em) -- +(1em,0) -- cycle;
%%    \end{tikzpicture}
%%}
%%% NEW PROBLEMS BOX
%
%\makeatletter
%\newlength{\fullwidthlength}
%\AtBeginDocument{\setlength{\fullwidthlength}{\@tufte@fullwidth}}
%\makeatother
%
%%%%%%% Formula with arrow
%\usepackage{mathtools}
%\usepackage{empheq}
%\definecolor{myblue}{rgb}{.8, .7, 1}
%\newcommand*\mybluebox[1]{%
%\colorbox{myblue}{\hspace{1em}#1\hspace{1em}}}
%
%\usepackage{amssymb}
%\usepackage[customcolors,markings]{hf-tikz}
%
%\tikzset{offset def/.style={
%    above left offset={0.015,0.3},
%    below right offset={-0.015,-0.12},
%  },
%  use color/.style={
%    offset def,
%    set fill color=white,
%    set border color=#1,
%  },
%}
%
%\newcommand{\annotate}[2][]{
%\tikz[remember picture,overlay]\draw[#1,use marker id] (0,0.4) -- +(90:0.7cm)
% node[anchor=north,text=black,align=center] {#2};
%}
%\usetikzlibrary{positioning}
%\tikzset{>=stealth}
%\newcommand{\tikzmark}[3][]{\tikz[overlay,remember picture,baseline] \node [anchor=base,#1](#2) {#3};}
%\usepackage[utf8]{inputenc}
%%%%%%% Formula with arrow
%
%%%%% FANCY COLOR TABLE
%\usepackage{tcolorbox}
%\usepackage{tabularx}
%\usepackage{array}
%\usepackage{colortbl}
%\tcbuselibrary{skins}
%\newcolumntype{Y}{>{\raggedleft\arraybackslash}X}
%
%\tcbset{tab1/.style={fonttitle=\bfseries\large,fontupper=\normalsize\sffamily,
%colback=yellow!10!white,colframe=red!75!black,colbacktitle=Salmon!40!white,
%coltitle=black,center title,freelance,frame code={
%\foreach \n in {north east,north west,south east,south west}
%{\path [fill=red!75!black] (interior.\n) circle (3mm); };},}}
%\tcbset{tab2/.style={enhanced,fonttitle=\bfseries,fontupper=\normalsize\sffamily,
%colback=yellow!10!white,colframe=red!50!black,colbacktitle=Salmon!40!white,
%coltitle=black,center title}}
%%%%% FANCY COLOR TABLE
%
%
%\usepackage{times}
%\colorlet{helpful}{lime!70}
%\colorlet{harmful}{red!30}
%\colorlet{internal}{yellow!20}
%\colorlet{external}{cyan!30}
%\colorlet{G}{helpful!50!internal}
%\colorlet{A}{harmful!50!internal}
%\colorlet{O}{helpful!50!external}
%\colorlet{T}{harmful!50!external}
%\newcommand{\texta}{Given}
%\newcommand{\textb}{Asking}
%\newcommand{\textcn}{\rotatebox[origin=c]{90}{\parbox[t]{2cm}{\centering Analyze\\ the Problem}}}
%\newcommand{\texts}{strength 1\par strength 2}
%\newcommand{\textw}{weakness 1\par weakness 2}
%\newcommand{\texto}{opportunity 1\par opportunity 2}
%\newcommand{\textt}{threat 1\par threat 2}
%\tcbset{swotbox/.style={size=normal, boxrule=0pt,
%            colback=#1, watermark text=#1, width=.4\linewidth-5mm},
%        header/.style={size=small, boxrule=0pt, width=.4\linewidth-5mm, colback=#1, valign=center, halign=center},
%        firstcol/.style={header=#1, width=1cm}
%        }
%        
%        
%        \tcbset{%
%  save height/.code={%
%    \tcbset{breakable}%
%    \providecommand{#1}{2cm}%
%    \def\tcb@split@start{%
%      \tcb@breakat@init%
%      \tcb@comp@h@page%
%      \def\tcb@ch{%
%        \tcbset{height=\tcb@h@page}%
%        \tcbdimto#1{#1+\tcb@h@page-\tcb@natheight}%
%        \immediate\write\@auxout{\string\gdef\string#1{#1}}%
%        \tcb@ch%
%      }%
%      \tcb@drawcolorbox@standalone%
%    }%
%  }%
%}
%
%%%%%%%%%% BACKGROUND IMAGE IN A BOX
%\newsavebox\MBox
%\newenvironment{Minipage}[1]
%  {\par\smallskip\begin{lrbox}{\MBox}\begin{minipage}{#1}}
%  {\end{minipage}\end{lrbox}%
%   \makebox(0,0){\put(-24,-19){%
%     \shadowimage[width=1.5\wd\MBox,height=2.1\ht\MBox]{paperpage}}}%
%%     \includegraphics[width=1.8\wd\MBox,height=2.1\ht\MBox]{paperpage}}}%
%   \usebox\MBox\par%
%  }
%\newenvironment{Minipage2}[1]
%  {\par\smallskip\begin{lrbox}{\MBox}\begin{minipage}{#1}}
%  {\end{minipage}\end{lrbox}%
%   \makebox(0,0){\put(-31,-19){%
%     \shadowimage[width=1.08\wd\MBox,height=2.4\ht\MBox]{paperpage}}}%
%%     \includegraphics[width=1.8\wd\MBox,height=2.1\ht\MBox]{paperpage}}}%
%   \usebox\MBox\par%
%  }
%%%%%%%%%% BACKGROUND IMAGE IN A BOX
%
%%%% SHADOW IN AN IMAGE
%\usetikzlibrary{shadows,calc}
%
%\def\shadowshift{3pt,-3pt}
%\def\shadowradius{6pt}
%
%\colorlet{innercolor}{black!60}
%\colorlet{outercolor}{gray!05}
%
%\newcommand\drawshadow[1]{
%    \begin{pgfonlayer}{shadow}
%        \shade[outercolor,inner color=innercolor,outer color=outercolor] ($(#1.south west)+(\shadowshift)+(\shadowradius/2,\shadowradius/2)$) circle (\shadowradius);
%        \shade[outercolor,inner color=innercolor,outer color=outercolor] ($(#1.north west)+(\shadowshift)+(\shadowradius/2,-\shadowradius/2)$) circle (\shadowradius);
%        \shade[outercolor,inner color=innercolor,outer color=outercolor] ($(#1.south east)+(\shadowshift)+(-\shadowradius/2,\shadowradius/2)$) circle (\shadowradius);
%        \shade[outercolor,inner color=innercolor,outer color=outercolor] ($(#1.north east)+(\shadowshift)+(-\shadowradius/2,-\shadowradius/2)$) circle (\shadowradius);
%        \shade[top color=innercolor,bottom color=outercolor] ($(#1.south west)+(\shadowshift)+(\shadowradius/2,-\shadowradius/2)$) rectangle ($(#1.south east)+(\shadowshift)+(-\shadowradius/2,\shadowradius/2)$);
%        \shade[left color=innercolor,right color=outercolor] ($(#1.south east)+(\shadowshift)+(-\shadowradius/2,\shadowradius/2)$) rectangle ($(#1.north east)+(\shadowshift)+(\shadowradius/2,-\shadowradius/2)$);
%        \shade[bottom color=innercolor,top color=outercolor] ($(#1.north west)+(\shadowshift)+(\shadowradius/2,-\shadowradius/2)$) rectangle ($(#1.north east)+(\shadowshift)+(-\shadowradius/2,\shadowradius/2)$);
%        \shade[outercolor,right color=innercolor,left color=outercolor] ($(#1.south west)+(\shadowshift)+(-\shadowradius/2,\shadowradius/2)$) rectangle ($(#1.north west)+(\shadowshift)+(\shadowradius/2,-\shadowradius/2)$);
%        \filldraw ($(#1.south west)+(\shadowshift)+(\shadowradius/2,\shadowradius/2)$) rectangle ($(#1.north east)+(\shadowshift)-(\shadowradius/2,\shadowradius/2)$);
%    \end{pgfonlayer}
%}
%
%\pgfdeclarelayer{shadow} 
%\pgfsetlayers{shadow,main}
%
%\newsavebox\mybox
%\newlength\mylen
%
%\newcommand\shadowimage[2][]{%
%\setbox0=\hbox{\includegraphics[#1]{#2}}
%\setlength\mylen{\wd0}
%\ifnum\mylen<\ht0
%\setlength\mylen{\ht0}
%\fi
%\divide \mylen by 120
%\def\shadowshift{\mylen,-\mylen}
%\def\shadowradius{\the\dimexpr\mylen+\mylen+\mylen\relax}
%\begin{tikzpicture}
%\node[anchor=south west,inner sep=0] (image) at (0,0) {\includegraphics[#1]{#2}};
%\drawshadow{image}
%\end{tikzpicture}}
%
%%%% SHADOW IN AN IMAGE
%%% TABULAR X ALIAS
%\newcolumntype{Y}{>{\centering\arraybackslash}X} % for centering tables in tabularx
%\def\tabularxcolumn#1{m{#1}}
%\newcommand{\heading}[1]{\bfseries\begin{tabular}{@{}c@{}} #1 \end{tabular}}
%%% TABULAR X ALIAS
%
%
%
%%%%%%%%%%TITLEPAGE
%\newcommand*{\rotrt}[1]{\rotatebox{90}{#1}} % Command to rotate right 90 degrees
%\newcommand*{\rotlft}[1]{\rotatebox{-90}{#1}} % Command to rotate left 90 degrees
%
%
%\newcommand*{\titleBC}{\begingroup % Create the command for including the title page in the document
% \thispagestyle{empty}
%\centering % Center all text
%
%\def\CP{\textit{\Huge Experiments for Allied health Chemistry}} % Title
%
%\settowidth{\unitlength}{\CP} % Set the width of the curly brackets to the width of the title
%{\color{LightGoldenrod}\resizebox*{\unitlength}{\baselineskip}{\rotrt{$\}$}}} \\[\baselineskip] % Print top curly bracket
%\textcolor{Sienna}{\CP} \\[\baselineskip] % Print title
%{\color{RosyBrown}\Large AN OPENTEXTBOOK} \\ % Tagline or further description
%{\color{Magenta}Version 1.0} \\ % Tagline or further description
%{\color{LightGoldenrod}\resizebox*{\unitlength}{\baselineskip}{\rotlft{$\}$}}} % Print bottom curly bracket
%
%\vfill % Whitespace between the title and the author name
% \includegraphics[width=0.8\textwidth]{coverfolder/cover}
% \vfill 
%%{\Large\textbf{Daniel Torres \& Luis Gonzalez-Urbina}}\\ % Author name
%{\textbf{B.M.C.C.}}\\ % Author name
%{\textbf{The City University of New York}}\\ % Author name
%\vfill % Whitespace between the author name and the publisher logo
%
% \includegraphics[width=0.1\textwidth]{coverfolder/cc}\\[0.5\baselineskip] % Publisher logo
%2022 % Year published
%
%%\newgeometry{textwidth=.58\paperwidth} % size of text column
%%\newgeometry{
%%  left=21.8mm, % left margin
%%  textwidth=125mm, % main text block
%%  marginparsep=8.2mm, % gutter between main text block and margin notes
%%  marginparwidth=49.4mm % width of margin notes
%%}
% \endgroup}%%%%%%%%%% TITLE PAGE
%
%
%
%
%
%
%
%
%
%
%%%%% PROBLEMS CHAPTER BOX
%\mdfdefinestyle{conclusion}{%
%innerleftmargin=0.3\linewidth,
%innerrightmargin=0.1cm,
%leftline=false,
%rightline=false,
%bottomline=false,
%backgroundcolor=white,
%linecolor=DarkGreen,
%linewidth=4pt,
%skipabove=\dimexpr\topskip+.5cm\relax,
%settings={\setlist[itemize]{label=\color{DarkGreen}\rule[.25ex]{1ex}{1ex}}%
%               \refstepcounter{section}},
%frametitleaboveskip=\dimexpr-1.5\ht\strutbox\relax,
%frametitle=%
% {%
%  \null\llap{%
%     \tikz[baseline=(current bounding box.east),outer sep=0pt,remember picture]
%      \node[anchor=west,
%             xshift=0cm,fill=DarkGreen,
%             rounded corners,
%             font=\color{white}\Huge\sffamily] (frametitle)%
%                                                                 {CHAPTER~\thechapter};
%               }%
%  }
%}
%
%\newmdenv[style=conclusion]{conclusion}
%
%
%\newcommand{\pol}[1]{\rlap{${}^{^{\color{red} \delta #1}}$}}
%\usepackage{fancybox}
%
%
%\usepackage{fontawesome}
%\def\Plus{\texttt{+}}
%\def\Minus{\texttt{-}}
%\newenvironment{myenv}[1]
%  {\mdfsetup{
%    frametitle={\colorbox{white}{\space#1\space}},
%    innertopmargin=0pt,
%    frametitleaboveskip=-\ht\strutbox,
%    frametitlealignment=\center
%    }
%  \begin{mdframed}%
%  }
%  {\end{mdframed}}
%
%%%%% BOXES DIAGRAMS  FOR CARBS
%\tikzset{
%     block/.style={rectangle, draw, fill=red!40, text width=4em,
%                   text centered, rounded corners, minimum height=3em},
%     arrow/.style={-{Stealth[]}}
%     }
%\usetikzlibrary{positioning,arrows.meta}
%%%%% BOXES DIAGRAMS  FOR CARBS
%
%
%
%\usepackage{enumitem}
%\SetLabelAlign{myright}{\hss\llap{$#1$}}
%\newlist{where}{description}{1}
%\setlist[where]{labelwidth=2cm,labelsep=1em,
%                        leftmargin=!,align=myright,font=\normalfont}
%
%
%
%\usepackage{smartdiagram}
%
%\usepackage{PTSansNarrow}
%\usepackage[T1]{fontenc}
%
%\usetikzlibrary{matrix}
%
%
%\usepackage{rotating}
%
%
%\def\parsecomma#1,#2\endparsecomma{\def\page@x{#1}\def\page@y{#2}}
%\tikzdeclarecoordinatesystem{page}{
%    \parsecomma#1\endparsecomma
%    \pgfpointanchor{current page}{north east}
%    \pgf@xc=\pgf@x%
%    \pgf@yc=\pgf@y%
%    \pgfpointanchor{current page}{south west}
%    \pgf@xb=\pgf@x%
%    \pgf@yb=\pgf@y%
%    \pgfmathparse{(\pgf@xc-\pgf@xb)/2.*\page@x+(\pgf@xc+\pgf@xb)/2.}
%    \expandafter\pgf@x\expandafter=\pgfmathresult pt
%    \pgfmathparse{(\pgf@yc-\pgf@yb)/2.*\page@y+(\pgf@yc+\pgf@yb)/2.}
%    \expandafter\pgf@y\expandafter=\pgfmathresult pt
%}
%\makeatother
%
%
%
%
%%%%% Add text in itemize
%\newcounter{saveenumerate}
%\usepackage{enumitem}
%\setlist[enumerate]{itemsep=0mm}
%\makeatletter
%\newcommand{\enumeratext}[1]{%
%\setcounter{saveenumerate}{\value{enum\romannumeral\the\@enumdepth}}
%\end{enumerate}
%#1
%\begin{enumerate}
%\setcounter{enum\romannumeral\the\@enumdepth}{\value{saveenumerate}}%
%}
%\makeatother
%
%
%
%
%
%
%\usetikzlibrary{arrows,positioning}
%
%%%%% COLOR ARROWS
%\usetikzlibrary{shapes.arrows}
%\makeatletter
%\def\createshadingfromlist#1#2#3{%
%  \pgfutil@tempcnta=0\relax
%  \pgfutil@for\pgf@tmp:={#3}\do{\advance\pgfutil@tempcnta by1}%
%  \ifnum\pgfutil@tempcnta=1\relax%
%    \edef\pgf@spec{color(0)=(#3);color(100)=(#3)}%
%  \else%
%    \pgfmathparse{50/(\pgfutil@tempcnta-1)}\let\pgf@step=\pgfmathresult%
%    %
%    \pgfutil@tempcntb=1\relax%
%    \pgfutil@for\pgf@tmp:={#3}\do{%
%      \ifnum\pgfutil@tempcntb=1\relax%
%        \edef\pgf@spec{color(0)=(\pgf@tmp);color(25)=(\pgf@tmp)}%
%      \else%
%        \ifnum\pgfutil@tempcntb<\pgfutil@tempcnta\relax%
%          \pgfmathparse{25+\pgf@step/4+(\pgfutil@tempcntb-1)*\pgf@step}%
%          \edef\pgf@spec{\pgf@spec;color(\pgfmathresult)=(\pgf@tmp)}%
%        \else%
%          \edef\pgf@spec{\pgf@spec;color(75)=(\pgf@tmp);color(100)=(\pgf@tmp)}%
%        \fi%
%      \fi%
%      \advance\pgfutil@tempcntb by1\relax%
%    }%
%  \fi%
%  \csname pgfdeclare#2shading\endcsname{#1}{100}\pgf@spec%
%}
%
%\createshadingfromlist{shading1}{horizontal}{red,yellow,green,cyan,blue}
%
%
%
%
%
%
%%ACID CONJGATE BASE
%\usetikzlibrary{calc,positioning}
%
%
%
%%%% INVERT COLOR MAP
%\def\customrevertcolormap#1{%
%    \pgfplotsarraycopy{pgfpl@cm@#1}\to{custom@COPY}%
%    \c@pgf@counta=0
%    \c@pgf@countb=\pgfplotsarraysizeof{custom@COPY}\relax
%    \c@pgf@countd=\c@pgf@countb
%    \advance\c@pgf@countd by-1 %
%    \pgfutil@loop
%    \ifnum\c@pgf@counta<\c@pgf@countb
%        \pgfplotsarrayselect{\c@pgf@counta}\of{custom@COPY}\to\pgfplots@loc@TMPa
%        \pgfplotsarrayletentry\c@pgf@countd\of{pgfpl@cm@#1}=\pgfplots@loc@TMPa
%        \advance\c@pgf@counta by1 %
%        \advance\c@pgf@countd by-1 %
%    \pgfutil@repeat
%%\pgfplots@colormap@showdebuginfofor{#1}%
%}%
%
%
%%Extra practice problems box
%\newtcolorbox[auto counter]{summary}[1][]{title={\bfseries Extra Practice Problems},enhanced,drop shadow={black!50!white},
%  coltitle=black,
%  top=0.3in,
%  attach boxed title to top right=
%  {xshift=0pt},
%  boxed title style={size=small,colback=pink},#1}
%%Extra practice problems box
%
%
%%BRACKET FOR CHEMFIG
%\newcommand\setpolymerdelim[2]{\def\delimleft{#1}\def\delimright{#2}}
%
%\def\makebraces(#1,#2)#3#4#5{%
%  \edef\delimhalfdim{\the\dimexpr(#1+#2)/2}%
%  \edef\delimvshift{\the\dimexpr(#1-#2)/2}%
%  \chemmove{
%    \node[at=(#4),yshift=(\delimvshift)]
%      {$
%       \left\delimleft
%         \vrule height\delimhalfdim depth\delimhalfdim width0pt
%       \right.
%      $};
%    \node[at=(#5),yshift=(\delimvshift)]
%      {$
%        \left.
%          \vrule height\delimhalfdim depth\delimhalfdim width0pt
%        \right\delimright_{\rlap{#3}}
%      $};
%  }%
%}
%
%
%\usepackage[weather]{ifsym}%
%%\usepackage{xymtexps}%PostScript Mode
%
%
%
%
%
%%%%%% DNA
%\usepackage{carbohydrates}
%\definesubmol\nobond{[,0.2,,,draw=none]}
%\definesubmol{adenine}{-[:72,,,,draw=none]\textcolor{blue}{N}(-[:-90]H)*5(-*6(-N=-N=(-NH_2)-=)--N=-)}
%\definesubmol{adenine2}{-[:72,0.4,,,draw=none]\textcolor{blue}{N}*5(-*6(-N=-N=(-NH_2)-=)--N=-)}
%\definesubmol{guanine}{-[:72,,,,draw=none]\textcolor{blue}{N}(-[:-90]H)*5(-*6(-N=(-[:-45]NH_2)-N(-[:45]H)-(=O)-=)--N=-)}
%\definesubmol{guanine2}{-[:72,0.4,,,draw=none]\textcolor{blue}{N}*5(-*6(-N=(-[:-45]NH_2)-N(-[:45]H)-(=O)-=)--N=-)}
%
%\definesubmol{uracil}{-[:90,,,,draw=none]\textcolor{blue}{N}(-[:-90]H)*6(-(=O)-N(-[:45]H)-(=O)-=-)}
%\definesubmol{uracil2}{-[:90,0.4,,,draw=none]\textcolor{blue}{N}*6(-(=O)-N(-[:45]H)-(=O)-=-)}
%
%\definesubmol{thymine}{-[:90,0.4,,,draw=none]\textcolor{blue}{N}(-[:-90]H)*6(-(=O)-N(-[:45]H)-(=O)-(-[:145]CH_3)=-)}
%\definesubmol{thymine2}{-[:90,0.4,,,draw=none]\textcolor{blue}{N}*6(-(=O)-N(-[:45]H)-(=O)-(-[:145]CH_3)=-)}
%
%\definesubmol{cytosine}{-[:90,,,,draw=none]\textcolor{blue}{N}(-[:-90]H)*6(-(=O)-N=(-NH_2)-=-)}
%\definesubmol{cytosine2}{-[:90,0.4,,,draw=none]\textcolor{blue}{N}*6(-(=O)-N=(-NH_2)-=-)}
%
%\definesubmol{ribose}{HO-[:90,0.5,2]?-[:60,1.0](-[:90,1.1]\textcolor{orange}{OH})-[:160,1.2]\textcolor{red}{O}-[:-160,1.2](-[:90,0.7](-[:180,0.9]H\textcolor{pink}{O}))-[:-60, 1.0]?(-[:-90,0.5]OH)}
%\definesubmol{deoxyribose}{-[:90,0.5,2,,,draw=none]?-[:60,1.0](-[:90,1.1]OH)-[:160,1.2]\textcolor{red}{O}-[:-160,1.2](-[:90,0.7](-[:180,0.9]H\textcolor{pink}{O}))-[:-60, 1.0]?(-[:-90,0.5]OH)}
%\definesubmol{phosphate}{-[6,.6]O-[6,1.5]P(-[4,.8]HO)(=[6,0.8]O)-[,1.5]O}
%\newcommand*\derivesubmol[4]{%
%    \saveexpandmode\saveexploremode\expandarg\exploregroups
%    \csname @\ifcat\relax\noexpand#2first\else second\fi oftwo\endcsname
%        {\expandafter\StrSubstitute\@car#2\@nil}
%        {\expandafter\StrSubstitute\csname CF@@#2\endcsname}
%    {\@empty#3}{\@empty#4}[\temp@]%
%    \csname @\ifcat\relax\noexpand#1first\else second\fi oftwo\endcsname
%        {\expandafter\let\@car#1\@nil}
%        {\expandafter\let\csname CF@@#1\endcsname}\temp@
%    \restoreexpandmode\restoreexploremode
%}
%\definesubmol{rt1}{-[2]rt1}
%\definesubmol{rt2}{-[6,.6]rt2}
%\definesubmol{rt3}{-[6,.6]rt3}
%\definesubmol{ribose2}{%
%    H\textcolor{pink}{O}-[:0,0.8]-[:-90,0.8]%
%        (%
%            -[:25,1.176]\textcolor{red}{O}%
%            -[:-25,1.176]%
%        )%
%    -[:-45,0.8]%
%        (%
%          (-[:-90,0.8]OH)-[0,,,,line width=1pt,shorten <=-.5pt,shorten >=-.5pt]%
%              (!{rt2})%
%          -[:45,0.8]%
%              (!{rt1})%
%         )%
%}
%\definesubmol{deoxyribose2}{%
%    H\textcolor{pink}{O}-[:0,0.8]-[:-90,0.8]%
%        (%
%            -[:25,1.176]\textcolor{red}{O}%
%            -[:-25,1.176]%
%        )%
%    -[:-45,0.8]%
%        (%
%          (-[:-90,0.8]OH)-[0,,,,line width=1pt,shorten <=-.5pt,shorten >=-.5pt]%
%              ()%
%          -[:45,0.8]%
%              (!{rt1})%
%         )%
%}
%%\derivesubmol{deoxyribose2}{ribose2}{(!{rt2})}{   }
%\definesubmol{deoxyribosephosphate}{%
%    P(-[:-90,0.8]OH)(-[:90,0.8]OH)(=[:180,0.8]O)-\textcolor{pink}{O}-[:0,0.8]-[:-90,0.8]%
%        (%
%            -[:25,1.176]\textcolor{red}{O}%
%            -[:-25,1.176]%
%        )%
%    -[:-45,0.8]%
%        (%
%          (-[:-90,0.8]OH)-[0,,,,,shorten <=-.5pt,shorten >=-.5pt]%
%          -[:45,0.8]%
%              (!{rt1})%
%         )%
%}
%\definesubmol{ribosephosphate}{%
%    P(-[:-90,0.8]OH)(-[:90,0.8]OH)(=[:180,0.8]O)-\textcolor{pink}{O}-[:0,0.8]-[:-90,0.8]%
%        (%
%            -[:25,1.176]\textcolor{red}{O}%
%            -[:-25,1.176]%
%        )%
%    -[:-45,0.8]%
%        (%
%          (-[:-90,0.8]OH)-[0,,,,,shorten <=-.5pt,shorten >=-.5pt]%
%          (-[:-90,0.8]OH)-[:45,0.8]%
%              (!{rt1})%
%         )%
%}
%%%%%% DNA
%
%
%
%%%%% AMINO ACIDS
%%chemical symbols
%\newcommand{\glycine}{\footnotesize\chemfig{C([:0]-C(=[:30]O)(-[:-30]O^{-}))([:180]-H_{3} N^{+})([:-90]-H)([:90]-H)}}
%%\newcommand{\proline}{\footnotesize\chemfig{C([:0]-C(=[:30]O)(-[:-30]O^{-}))([:-90]-H)*5([:135]-CH_2-CH_2-H_2C-H_{2} N^{+}-)}}
%\newcommand{\alanine}{\footnotesize\chemfig{C([:0]-C(=[:30]O)(-[:-30]O^{-}))([:180]-H_{3}
%    N^{+})([:90]-CH_{3})([:-90]-H)}}
%\newcommand{\valin}{\footnotesize\chemfig{C([:0]-C(=[:30]O)(-[:-30]O^{-}))([:180]-H_{3} N^{+})([:90]-CH([:120]-CH_{3})([:60]-CH_{3}))([:-90]-H)}}
%\newcommand{\leucin}{\footnotesize\chemfig{C([:0]-C(=[:30]O)(-[:-30]O^{-}))([:180]-H_{3}
%    N^{+})([:-90]-H)([:90]-CH_{2}([:90]-CH([:120]-CH_{3})([:60]-CH_{3})))}}
%\newcommand{\isoleucin}{\footnotesize\chemfig{C([:0]-C(=[:30]O)(-[:-30]O^{-}))([:180]-H_{3}
%    N^{+})([:-90]-H)([:90]-CH ([:90]-CH_{2}([:90]-CH_3))
%    ([:180]-H_{3}C))}}
%\newcommand{\methionin}{\footnotesize\chemfig{C([:0]-C(=[:30]O)(-[:-30]O^{-}))([:180]-H_{3}
%    N^{+})([:-90]-H)([:90]-CH_{2}([:90]-CH_{2}([:90]-S ([:90]-CH_{3}))))}}
%\newcommand{\phenylalanin}{\footnotesize\chemfig{C([:0]-C(=[:30]O)(-[:-30]O^{-}))([:180]-H_{3}
%    N^{+})([:-90]-H)([:90]-CH_{2}([:90]-*6(-=-=-=)))}}
%\newcommand{\tryptophan}{\footnotesize\chemfig{C([:0]-C(=[:30]O)(-[:-30]O^{-}))([:180]-H_{3}
%    N^{+})([:-90]-H)([:90]-CH_{2}([:90]-([:60]*6(?=-NH-*6([:120]-=-=-)=?))))}}
%\newcommand{\serin}{\footnotesize\chemfig{C([:0]-C(=[:30]O)(-[:-30]O^{-}))([:180]-H_{3}
%    N^{+})([:-90]-H)([:90]-CH_{2}([:90]-OH))}}
%\newcommand{\threonin}{\footnotesize\chemfig{C([:0]-C(=[:30]O)(-[:-30]O^{-}))([:180]-H_{3}
%    N^{+})([:-90]-H)([:90]-CH_{2}([:90]-CH([:120]-OH)([:60]-CH_{3})))}}
%\newcommand{\cystein}{\footnotesize\chemfig{C([:0]-C(=[:30]O)(-[:-30]O^{-}))([:180]-H_{3}
%    N^{+})([:-90]-H)([:90]-CH_{2}([:90]-CH_{2}([:90]-SH)))}}
%\newcommand{\tyrosin}{\footnotesize\chemfig{C([:0]-C(=[:30]O)(-[:-30]O^{-}))([:180]-H_{3}
%    N^{+})([:-90]-H)([:90]-CH_{2}([:90]-*6(-=-([:90]-OH)=-=)))}}
%\newcommand{\asparagin}{\footnotesize\chemfig{C([:0]-C(=[:30]O)(-[:-30]O^{-}))([:180]-H_{3}
%    N^{+})([:-90]-H)([:90]-CH_{2}([:90]-C([:120]-NH_{2})([:60]=O)))}}
%\newcommand{\glutamin}{\footnotesize\chemfig{C([:0]-C(=[:30]O)(-[:-30]O^{-}))([:180]-H_{3}
%    N^{+})([:-90]-H)([:90]-CH_{2}([:90]-CH_{2}([:90]-C([:120]-NH_{2})([:60]=O))))}}
%\newcommand{\asparaginacid}{\footnotesize\chemfig{C([:0]-C(=[:30]O)(-[:-30]O^{-}))([:180]-H_{3}
%    N^{+})([:-90]-H)([:90]-CH_{2}([:90]-C([:120]-O^{-})([:60]=O)))}}
%\newcommand{\glutaminacid}{\footnotesize\chemfig{C([:0]-C(=[:30]O)(-[:-30]O^{-}))([:180]-H_{3}
%    N^{+})([:-90]-H)([:90]-CH_{2}([:90]-CH_{2}([:90]-C([:120]-O^{-})([:60]=O))))}}
%\newcommand{\lysin}{\footnotesize\chemfig{C([:0]-C(=[:30]O)(-[:-30]O^{-}))([:180]-H_{3}
%    N^{+})([:-90]-H)([:90]-CH_{2}([:90]-CH_{2}([:90]-CH_{2}([:90]-CH_{2}([:90]-NH_{3}^{+})))))}}
%\newcommand{\arginin}{\footnotesize\chemfig{C([:0]-C(=[:30]O)(-[:-30]O^{-}))([:180]-H_{3}
%    N^{+})([:-90]-H)([:90]-CH_{2}([:90]-CH_{2}([:90]-CH_{2}([:90]-NH([:90]-C([:90]-NH_{2})([:0]=NH_2^{+}))))))}}
%%\newcommand{\histidin}{\footnotesize\chemfig{C([:0]-C(=[:30]O)(-[:-30]O^{-}))([:180]-H_{3}N^{+})([:-90]-H)([:90]-CH_{2}([:90]-([:60]*5(-NH-=NH^{+}-=))))}}
%
%\newcommand{\proline}{\footnotesize
%\chemfig{
%  H-[2]
%  C
%  (-C(=[:45]O)(-[:-45]O|^{-}))
% *5([:135]-C|H_2-\chemabove{C}{H\mathrlap{_2}}-H_2|C-[,,,2]H_{2}|\chemabove{N}{\quad\scriptstyle+}-)
%}
%}
%
%\newcommand{\histidin}{\footnotesize
%\chemfig{
%  H_{3}\chemabove{N}{\quad\scriptstyle+}-
%  C
%   ([:-90]-H)
%   ([:90]
%     -CH_{2}-
%     ([:90]*5(-NH-[,,1]=\chemabove{\chemabove{N}{H}}{\hskip1.5em\scriptstyle+}-=))
%   )
%  -C(=[:45]O)(-[:-45]O|^{-})
%}
%}
%
%
%
%% BRACET LIPID
%\usetikzlibrary{decorations.pathreplacing}
%\tikzset{brace/.style = { decorate, decoration={brace, amplitude=5pt} },}
%% BRACET LIPID
%
%%%%LIPIDS DISGRAMS
%\@namedef{color@1}{red!40}
%\@namedef{color@2}{green!40}   
%\@namedef{color@3}{blue!40} 
%\@namedef{color@4}{cyan!40}  
%\@namedef{color@5}{magenta!40} 
%\@namedef{color@6}{yellow!40}    
%
%\newcommand{\graphitemize}[2]{%
%\begin{tikzpicture}[every node/.style={align=center}]  
%  \node[minimum size=4cm,circle,fill=gray!40,font=\Large,outer sep =.25cm,inner sep=.5cm](ce){#1};  
%\foreach \gritem [count=\xi] in {#2}  {\global\let\maxgritem\xi}  
%
%\foreach \gritem [count=\xi] in {#2}
%{% 
%\pgfmathtruncatemacro{\angle}{360/\maxgritem*\xi}
%\edef\col{\@nameuse{color@\xi}}
%\node[circle,
%     ultra thick,
%     draw=white,
%     fill opacity=.5,
%     fill=\col,outer sep=0.25cm,        
%     minimum size=2cm] (satellite-\xi) at (\angle:5cm) {\gritem };
%     \draw[line width=0.5cm,->,\col] (ce) -- (satellite-\xi);
%     }%
%\end{tikzpicture}  
%}%
%
%
%
%
%%%%% CHECKBOX
%\usepackage{amssymb,tikz,pdftexcmds,xparse}
%\tikzset{box/.style={
%    minimum size=0.225cm,
%    inner sep=0pt,
%    draw,
%  },  
%  insert mark/.style={
%    append after command={%
%         node[inner sep=0pt,#1]
%           at (\tikzlastnode.center){$\checkmark$}
%     }     
%  },
%  insert bad mark/.style={
%    append after command={%
%         [shorten <=\pgflinewidth,shorten >=\pgflinewidth]
%         (\tikzlastnode.north west)edge[#1](\tikzlastnode.south east)
%         (\tikzlastnode.south west)edge[#1](\tikzlastnode.north east)
%     }     
%  },
%}
%\makeatletter
%\NewDocumentCommand{\tikzcheckmark}{O{} m}{%
%  \ifnum\pdf@strcmp{#2}{mark}=\z@%
%    \tikz[baseline=-0.5ex]\node[box,insert mark={#1},#1]{};%
%  \fi%
%  \ifnum\pdf@strcmp{#2}{bad mark}=\z@%
%    \tikz[baseline=-0.5ex]\node[box,insert bad mark={#1},#1]{};%
%  \fi%
%  \ifnum\pdf@strcmp{#2}{no mark}=\z@%
%    \tikz[baseline=-0.5ex]\node[box,#1]{};%
%  \fi%
%}
%\makeatother
%%%%%CHECKBOX
%
%
%
%\definecolor{mybrown}{RGB}{128,64,0}
%\usepackage{everypage}%%% SIDEBAR ON PAGES
%\usepackage{scrlayer}
%\DeclareNewLayer[
%  background,
%  outermargin,
%  addhoffset=\marginparsep,
%  contents={\color{gray!50}\rule{\layerwidth}{\layerheight}}
%]{sidebar.bg}
%\makeatletter
%\DeclareNewLayer[
%  clone=sidebar.bg,
%  addhoffset=1em,
%  addwidth=-1em,
%  contents={\parbox[b][\layerheight]{\layerwidth}{\raggedright\@sidebar}}
%]{sidebar.txt}
%\newcommand\@sidebar{}
%\newcommand\sidebar[1]{\renewcommand\@sidebar{#1}}
%\makeatother
%\DeclareNewPageStyleByLayers{sidebar}{sidebar.bg,sidebar.txt}%%% SIDEBAR ON PAGES
%\usepackage{makecell}
%
%
%
%
%\definecolor{TableBorder}{rgb}{0.0078,0.6824,0.7882}
%
%\definecolor{TableEven}{rgb}{0.8000,0.9216,0.9490}
%
%\definecolor{TableOdd}{rgb}{1,1,1}
%
%\DeclareDocumentCommand\newstep{o}{%
%\item\IfNoValueTF{#1}{}{#1 \textendash\xspace}}
%\newlist{steps}{enumerate}{1}
%\setlist[steps]{label={} \textit{Step \arabic*:},leftmargin=*}
%
%\newwrite\tempfile
%\makeatletter
%\DeclareDocumentCommand \sidenote {m O{Pbardefault} O{Plightdefault}} {%
%\stepcounter{localsidenotecounter}%
%\hyperlink{sidenote\arabic{localsidenotecounter}}{${}^\textrm{\footnotesize\arabic{localsidenotecounter}}$}%
%\immediate\write\tempfile%
%{%
%\noexpand\Pnotebox{#2}{#3}{%
%\noexpand\footnotesize \noexpand{\noexpand\hypertarget{sidenote\arabic{localsidenotecounter}}{${}^{\arabic{localsidenotecounter}}$} \unexpanded{#1 }}%
%}}}
%\makeatother
%\usepackage{array,multirow,graphicx}
%
%
%%Numbers in circle
%\newcommand{\mycircled}[1]{%
% \begin{tikzpicture}[baseline={(char.base)}]
%   \node (char) {#1};
%   \node[draw,circle,minimum size=16pt,inner sep=0pt,overlay] (char.center){};
% \end{tikzpicture}
% }
%%Numbers in circle
%
%
%
%%nucleotides
%
%\derivesubmol{deoxyadenosine}{deoxyribose2}{!{rt1}}{-[2,1.0]!{adenine2}}
%\derivesubmol{deoxythymidine}{deoxyribose2}{!{rt1}}{-[2,1.0]!{thymine2}}
%\derivesubmol{deoxyuridine}{deoxyribose2}{!{rt1}}{-[2,1.0]!{uracil2}}
%\derivesubmol{deoxyguaninephosphate}{deoxyribosephosphate}{!{rt1}}{-[2,1.0]!{guanine2}}
%\derivesubmol{deoxyuridinephosphate}{deoxyribosephosphate}{!{rt1}}{-[2,1.0]!{uracil2}}
%\derivesubmol{thyminephosphate}{ribosephosphate}{!{rt1}}{-[2,1.0]!{thymine2}}
%\derivesubmol{cytosinephosphate}{ribosephosphate}{!{rt1}}{-[2,1.0]!{cytosine2}}
%\derivesubmol{deoxyadenosinephosphate}{deoxyribosephosphate}{!{rt1}}{-[2,1.0]!{adenine2}}
%
%
%
%%% CIRCLE LEWIS
%\newcommand*\circleatom[1]{\tikz\node[circle, fill=green!30]{\printatom{#1}};}
%%% CIRCLE LEWIS
%
%%%% ORBITAL DIAGRAMS
%\usepackage{tcolorbox}% no need in this mwe
%\usepackage[font=small]{caption}% no need in this mwe
%\usepackage{float}% no need in this mwe
%\usepackage{floatflt}% no need in this mwe
%\usepackage{tikz}
%\usepackage{tikzorbital}
%\usetikzlibrary{fadings,patterns,backgrounds,fit,arrows}
%\pgfdeclarelayer{backbackground}
%\pgfsetlayers{backbackground,background,main}
%\usepackage{amssymb} % for harpoons
%\newcommand{\electron}[2]{{%
%        \newcommand*\up{\fbox{$\mathord\upharpoonleft\phantom{\downharpoonright}$}}%
%        \newcommand*\dwn{\fbox{$\mathord\downharpoonleft\phantom{\upharpoonright}$}}%
%        \newcommand*\updwn{\fbox{$\upharpoonleft\downharpoonright$}}%
%        \newcommand*\emp{\fbox{$\phantom{\downharpoonright}\phantom{\downharpoonright}$}}%
%        \setlength\tabcolsep{0pt}% remove extra horizontal space from tabular
%        \begin{tabular}{c}#2\\[2pt]#1\end{tabular}}}%
%%%% ORBITAL DIAGRAMS
%
%%%%%EXAMPLE BOX
%\definecolor{mybrown}{RGB}{128,64,0}
%\mdfdefinestyle{mystyle}{%
%linecolor=black,outerlinewidth=1pt,%
%frametitlerule=true,frametitlefont=\sffamily\bfseries\color{white},%
%frametitlerulewidth=1pt,frametitlerulecolor=black,%
%frametitlebackgroundcolor=black,
%backgroundcolor=mybrown!05,
%innertopmargin=\topskip,
%roundcorner=5pt
%}
%\mdtheorem[style=mystyle]{example}{Example}
%
%\gdef\Sepline{%
%  \par\noindent\makebox[\linewidth][l]{%
%  \hspace*{-\mdflength{innerleftmargin}}%
%   \tikz\draw[thick,dashed,gray!60] (0,0) --%
%        (\textwidth+\the\mdflength{innerleftmargin}+\the\mdflength{innerrightmargin},0);
%  }\par\nobreak}
%%%%%EXAMPLE BOX
%
%\definecolor{tableShade}{gray}{0.9}
%
%
%\definecolor{bidentitlebg}{RGB}{158,59,255}%YOUR TUBE
%\newtcolorbox{ridentidad}[1][]{
%  enhanced,
%  frame code={
%    \fill[draw=white,top color=white,bottom color=blue!60]
%      ([xshift=-20pt]title.north west) --
%      (title.south west) --
%      (title.south east) --
%      ([xshift=20pt]title.north east) -- cycle;
%    \draw[blue,line width=0.4mm,rounded corners]
%      (frame.north west) --
%      (frame.south west) --
%      ([xshift=-20pt]title.north west) --
%      (title.south west) --
%      (title.south east) --
%      ([xshift=20pt]title.north east) --
%      (frame.south east) --
%      (frame.north east);
%  },
%  coltitle=blue!75!black,
%  colback=white,
%  attach boxed title to bottom center,
%  boxed title style={empty},
%  fonttitle=\bfseries\sffamily,
%  title=\strut Visual,
%  #1,
%}%YOUR TUBE
%
%
%\usetikzlibrary{quotes,arrows.meta}
%
%
%
%
%
%
%
%
%%%%%% COLOR TABLE
%\usepackage{environ}
%\usepackage{longtable}
%\newcommand*\rot{\rotatebox{90}}
%\usepackage{chemfig}
%\def\Z{  0 & 0 & 12 & Working with leaflets & Claudio Fiandrino\\}
%\def\ZZ{\Z\Z\Z\Z}\def\ZZZ{\ZZ\ZZ\ZZ\ZZ}\def\ZZZZ{\ZZZ\ZZZ\ZZZ\ZZZ}
%
%% some stuff common to short and long tables
%\definecolor{tablecolor}{named}{ForestGreen}
%\newcommand{\header}[1]{%
%  \multicolumn{1}{c}{%
%    \cellcolor{tablecolor}\color{white}\bfseries#1}}
%\newcommand{\rndtablesetup}{%
%  \addtolength{\extrarowheight}{1ex}%
%  \rowcolors{2}{tablecolor!20}{tablecolor!40}%
%  \sffamily}
%
%% styles for the clipping/internal nodes
%\tikzset{
%  myclip/.style={rounded corners=1ex},
%  mynode/.style={anchor=south west,inner sep=0pt,outer sep=0pt},
%  mybord/.style={myclip,very thick,tablecolor},
%}
%
%% the short table (provided)
%\newsavebox{\tablebox}
%\NewEnviron{rndtable}[1]{%
%  \rndtablesetup%
%  \savebox{\tablebox}{%
%    \begin{tabular}{#1}%
%      \BODY%
%    \end{tabular}}%
%  \begin{tikzpicture}
%    \clip[myclip] (0,-\dp\tablebox) -- (\wd\tablebox,-\dp\tablebox) --
%      (\wd\tablebox,\ht\tablebox) -- (0,\ht\tablebox) -- cycle;
%    \node[mynode] at (0,-\dp\tablebox) {\usebox{\tablebox}};
%    \draw[mybord] (0,-\dp\tablebox) -- (\wd\tablebox,-\dp\tablebox) --
%      (\wd\tablebox,\ht\tablebox) -- (0,\ht\tablebox) -- cycle;
%  \end{tikzpicture}}
%
%% round top
%\newcommand{\roundtop}[1]{%
%  \begin{tikzpicture}
%    \clip[myclip] (0,0) |- (\wd#1,\ht#1+\dp#1) -- (\wd#1,0);
%    \draw[mybord] (0,0) |- (\wd#1,\ht#1+\dp#1) -- (\wd#1,0);
%    \node[mynode] at (0,0) {\box#1};
%  \end{tikzpicture}}
%
%% round bottom
%\newcommand{\roundbot}[1]{%
%  \begin{tikzpicture}
%    \clip[myclip] (0,0) |- (\wd#1,-\ht#1+\dp#1) -- (\wd#1,0);
%    \draw[mybord] (0,0) |- (\wd#1,-\ht#1+\dp#1) -- (\wd#1,0);
%    \node[mynode] at (0,0) {\box#1};
%  \end{tikzpicture}}
%
%% frame midsections
%\newcommand{\framemid}[1]{%
%  \begin{tikzpicture}
%    \draw[mybord] (0,0) -- (0,\ht#1+\dp#1) (\wd#1,0) -- (\wd#1,\ht#1+\dp#1);
%    \node[mynode] at (0,0) {\box#1};
%  \end{tikzpicture}}
%
%
%
%
%\usepackage{tcolorbox}% no need in this mwe
%\usepackage[font=small]{caption}% no need in this mwe
%\usepackage{float}% no need in this mwe
%\usepackage{floatflt}% no need in this mwe
%\usepackage{tikz}
%\usepackage{tikzorbital}
%\usetikzlibrary{fadings,patterns,backgrounds,fit,arrows}
%\pgfdeclarelayer{backbackground}
%\pgfsetlayers{backbackground,background,main}
%
%
%
%
%\usepackage{adjustbox}
%
%
%
% 
%
%
%
%
%\begin{document} 
%\thispagestyle{empty}
%\titleBC\newpage
%\clearpage\thispagestyle{empty}\mbox{}\clearpage
%%\section{Letter from the authors}\subfile{letters/letter2.tex}
%\clearpage\thispagestyle{empty}\mbox{}\clearpage
%  \clearpage\thispagestyle{empty}\mbox{}\clearpage   % COMMENT FOR CH202
%{ \hypersetup{hidelinks} \tableofcontents}
%%\clearpage\thispagestyle{empty}\mbox{}\clearpage
%%\newpage\section{Safety }\subfile{safety/safety.tex}
%%\newpage\section{Class policies}\subfile{safety/policies.tex}
%
%%\newpage\section{Hazard Symbols}\subfile{hazard/hazard.tex}
% \clearpage\mbox{}\clearpage
% \clearpage\mbox{}\clearpage
%
%
%
%%\subfile{./cover/title.tex}
%%\subfile{./periodictable/periodictable.tex}
%%\tableofcontents
%
%%\subfile{./chapter1/chapter0Exp.tex}
%%\subfile{./chapter1/chapter01Exp.tex}
%
%%\subfile{./chapter1/chapter1Exp.tex}
%%\subfile{./chapter2/chapter2Exp.tex}
%%\subfile{./chapter3/chapter3Exp.tex}
%%\subfile{./chapter4/chapter4Exp.tex}
%%\subfile{./chapter5/chapter5Exp.tex}
%%\subfile{./chapter6/chapter6Exp.tex}
%%\subfile{./chapter8/chapter8Exp.tex}
%%\subfile{./chapter9/chapter9Exp.tex}
%%\subfile{./chapter10/chapter10Exp.tex}
%%\subfile{./chapter11/chapter11Exp.tex}
%%\subfile{./chapter12/chapter12Exp.tex}
%%\subfile{./chapter13/chapter13Exp.tex}
%
%
%
%\clearpage\thispagestyle{empty}\mbox{}\clearpage
%
%
%
%
%
%
%
%\end{document} 






\documentclass{book}


\usepackage[svgnames,table,x11names]{xcolor} % Required to specify font color
\usepackage{lscape}
\usepackage{graphicx}% http://ctan.org/pkg/graphicx
\usepackage{array}% http://ctan.org/pkg/array
\newcommand*{\plogo}{\fbox{$\mathcal{PL}$}} % Generic publisher logo
\usepackage{lmodern}
\usepackage{graphicx} % Required for box manipulation
\usepackage{incgraph}

\usepackage{diagbox}
\usepackage[clock]{ifsym}
\usepackage{placeins}
\usepackage{tikz}
\usepackage{paralist}

\usetikzlibrary{plotmarks}
\begin{filecontents}{calor}
#time 	T
1 10
2  20
3 28
4 29
5 30
6 29.5
7 29
8 28.5
9  27.5
10 27
\end{filecontents}


\usepackage{float}

%\usepackage[many]{tcolorbox}




%----------------------------------------------------------------------------------------
%	TITLE PAGE
%----------------------------------------------------------------------------------------

\newcommand*{\rotrt}[1]{\rotatebox{90}{#1}} % Command to rotate right 90 degrees
\newcommand*{\rotlft}[1]{\rotatebox{-90}{#1}} % Command to rotate left 90 degrees

\newcommand*{\titleBC}{\begingroup % Create the command for including the title page in the document
%\includegraphics[width=1\textwidth]{cover0}
%\newpage
%\incgraph[documentpaper,
%  overlay={\node[red] at (page.center) {\Huge Picture sized to paper};}]
%  [width=\paperwidth,height=\paperheight]{cover0}
%\newpage
\centering % Center all text

\def\CP{\textit{\Huge Chemistry Experiments for the Health sciences}} % Title

\settowidth{\unitlength}{\CP} % Set the width of the curly brackets to the width of the title
{\color{LightGoldenrod}\resizebox*{\unitlength}{\baselineskip}{\rotrt{$\}$}}} \\[\baselineskip] % Print top curly bracket
\textcolor{Sienna}{\CP} \\[\baselineskip] % Print title
{\color{RosyBrown}\Large AN OPENTEXTBOOK} \\ % Tagline or further description
{\color{Magenta}Version 0.5 - \today } \\ % Tagline or further description
{\color{LightGoldenrod}\resizebox*{\unitlength}{\baselineskip}{\rotlft{$\}$}}} % Print bottom curly bracket

\vfill % Whitespace between the title and the author name
 \includegraphics[width=0.8\textwidth]{coverfolder/cover}
 \vfill 
%{\Large\textbf{Daniel Torres \& Luis Gonz\'alez-Urbina}}\\ % Author name
{\textbf{Science Department}}\\ % Author name
{\textbf{B.M.C.C.}}\\ % Author name
{\textbf{The City University of New York}}\\ % Author name
%{\textbf{www.materialsstemlab.nyc/education}}\\ % Author name
\vfill % Whitespace between the author name and the publisher logo

 \includegraphics[width=0.1\textwidth]{coverfolder/cc}\\[0.5\baselineskip] % Publisher logo

%\hspace{8cm} \includegraphics[width=0.2\textwidth, angle=45]{stamp}\\[0.5\baselineskip] % Publisher logo

\endgroup}








%----------------------------------------------------------------------------------------
%	BLANK DOCUMENT
%----------------------------------------------------------------------------------------
\usepackage{subfiles}
\usepackage[english]{babel}
\usepackage{import}
 \usepackage{subfiles}
%%%NEW
\usepackage{graphicx,lastpage}
\hyphenpenalty 10000
\flushbottom
\usepackage[normalem]{ulem}
\renewcommand\ULthickness{2pt}   %%---> For changing thickness of underline
\setlength\ULdepth{1.5ex}%\maxdimen ---> For changing depth of underline
\renewcommand{\baselinestretch}{1}
\pagestyle{empty}
\usepackage{layouts}
\usepackage[margin=0.7in,headsep=10sp]{geometry}
%\usepackage[margin=0.7in,headheight=65pt,headsep=1cm]{geometry}



\newcommand{\continuedmessage}{%
\ifcontinuation{\footnotesize Question \ContinuedQuestion\ continues\ldots}{}%
}
\usepackage[framemethod=tikz]{mdframed}
\usepackage{lipsum}
\mdfdefinestyle{mystyle}{%
linecolor=green!40!black,outerlinewidth=1pt,%
frametitlerule=true,frametitlefont=\sffamily\bfseries\color{white},%
frametitlerulewidth=1pt,frametitlerulecolor=green!40!black,%
frametitlebackgroundcolor=green!40!black,
backgroundcolor=green!5,
innertopmargin=\topskip,
roundcorner=1pt
}
\newmdenv[style=mystyle]{exa}
\newenvironment{example}[1]
  {\begin{exa}[frametitle=#1]}
  {\end{exa}}

\usepackage{cleveref}
\crefname{figure}{figure}{figures}
\crefname{question}{question}{questions}
\usepackage[most]{tcolorbox}
\usepackage{multicol}
\usepackage{enumitem}
\usepackage{amsfonts}
\setlist[itemize,1]{label=$\times$}
\setlist[itemize,2]{label=$\checkmark$}
\setlist[itemize,3]{label=$\diamond$}
\setlist[itemize,4]{label=$\bullet$}
\usepackage{tikz, adjustbox}
\usetikzlibrary{shapes,snakes}
\usepackage{amsmath,amssymb}
\usepackage{fourier}
\usepackage{wrapfig}
\usepackage{hyperref}
\usepackage{enumitem}
\usepackage{xspace}
\usepackage{xparse}
\DeclareDocumentCommand\newstep{o}{%
\item\IfNoValueTF{#1}{}{#1 \textendash\xspace}}
\newlist{steps}{enumerate}{1}
\setlist[steps]{label={\huge $\square$} \textit{Step \arabic*:},leftmargin=*}
\usepackage{marvosym}\usepackage{pifont}
\usepackage{tikz}
\usepackage{xparse}
\usetikzlibrary{tikzmark,decorations.pathreplacing,calc,decorations.markings}
\definecolor{colori}{RGB}{166,35,41}
\definecolor{colorii}{RGB}{248,219,162}
\usetikzlibrary{chains, positioning}
\usetikzlibrary{shadows,arrows.meta,positioning,backgrounds,fit,chains,scopes}

% Define block styles
\tikzset{%
  materia/.style={draw, fill=blue!20, text width=6.0em, text centered, minimum height=1.5em,drop shadow},
  etape/.style={materia, text width=8em, minimum width=10em, minimum height=3em, rounded corners, drop shadow},
  linepart/.style={draw, thick, color=black!50, -LaTeX, dashed},
  line/.style={draw, thick, color=black!50, -LaTeX},
  ur/.style={draw, text centered, minimum height=0.01em},
  back group/.style={fill=yellow!20,rounded corners, draw=black!50, dashed, inner xsep=15pt, inner ysep=10pt},
}


\NewDocumentCommand\MyArrow{O{0pt}mmmO{out=150,in=210}}
{%
\begin{tikzpicture}[overlay, remember picture]
  \draw [->,thick,line width=4pt,#4]
    ( $ ({pic cs:#3}|-{pic cs:#2}) + (-#1,1.3ex) $ ) to[#5]  
    ( $ (pic cs:#3) + (-#1,0) $ );
\end{tikzpicture}%
}
\usepackage[utf8]{inputenc}
 
\setlength{\arrayrulewidth}{1mm}
\setlength{\tabcolsep}{18pt}
\renewcommand{\arraystretch}{1.5}
\usepackage{array}\usepackage{multirow}\usepackage{gensymb}
\usepackage{graphicx}
\usepackage{caption}
\usepackage{subcaption}\usepackage{float}
\usepackage{tikzpagenodes}
\usepackage{lipsum}
\usepackage[some]{background}

\definecolor{notepadrule}{RGB}{217,244,244}
\backgroundsetup{
contents={%   
  \begin{tikzpicture}
    \foreach \fila in {0,...,52}
    {
      \draw [line width=1pt,color=notepadrule] 
      (current page.west|-0,-\fila*12pt) -- ++(\paperwidth,0);
    }
    \draw[overlay,red!70!black,line width=1pt]
      ([xshift=-1pt]current page text area.west|-current page.north) --  
      ([xshift=-1pt]current page text area.west|-current page.south);
  \end{tikzpicture}%
},
scale=1,
angle=0,
opacity=1
}
\usepackage{mathtools}
\usepackage[version=4]{mhchem}
\usepackage{tikz}
\newcommand{\mycircled}[1]{%
 \begin{tikzpicture}[baseline={(char.base)}]
   \node (char) {#1};
   \node[draw,circle,minimum size=16pt,inner sep=0pt,overlay] (char.center){};
 \end{tikzpicture}
 }
 
 
 \usepackage[newcommands]{ragged2e}
\usepackage{newfile}
\usepackage{color}
\usepackage{xparse}
\usepackage{graphics}
\usepackage{multirow}
\usepackage{ifthen}

% Settings
\colorlet{Pbardefault}{yellow!70} % Default color for the colored bar left to the notes
\colorlet{Plightdefault}{yellow!10}% Default color for background of the note
\def\Pmin{0.5} % ~minimal amount main text
\def\Pmax{0.95} % ~maximal amount main text / \Psf
\def\Psf{1.14} % 


\newcounter{localsidenotecounter}
\newlength{\currentparskip}
\setlength{\currentparskip}{\parskip}


% The formatting of the note
\newdimen\Pparsep
\Pparsep=0.2em

\newdimen\Pcolorwidth
\Pcolorwidth=2mm


% The paragraph command
\makeatletter
\def\sidenotepi#1\par%
{
%
\immediate\openout\tempfile=sidenotes.tex %
%
\twopara%
{
#1   \immediate\closeout\tempfile 
}
{\input{sidenotes.tex}}
}

\makeatother




% A command for finding the optimal text widths
\makeatletter

\long\def\twopara#1#2{%

\newdimen\columnsep
\columnsep=\dimexpr \Pparsep + \Pcolorwidth

 \def\@elt##1{\global\value{##1}\the\value{##1}\relax}% 
\edef\TY@ckpt{\cl@@ckpt}%

\@tempdima\dimexpr\textwidth-\columnsep\relax
{\hbadness\@M\raggedright\hsize.5\@tempdima \@tempdima\hsize
\setbox\z@\vbox{{#1\endgraf}}%
\setbox\tw@\vbox{{#2\endgraf}}%
\Gscale@div\tmp{\ht\z@}{\dimexpr\ht\z@+\ht\tw@\relax}%
\global\let\xtmp\tmp}%
\dimen@ \xtmp\@tempdima
\ifdim\dimen@<\Pmin\@tempdima \dimen@\Pmin\@tempdima\fi
\ifdim\dimen@>\Pmax\@tempdima \dimen@\Pmax\@tempdima\fi
\dimen@ii\dimexpr\@tempdima-\dimen@\relax 

\TY@ckpt


\ifdim\dimen@ii<20pt
  #1
\else
  \ifdim \Psf\dimen@ii>\dimexpr\textwidth-\Pmin\textwidth
    #1
    \let\Pnotebox\Pnotebox@
    \dimen@ii = \dimexpr \textwidth-4mm
      \setlength{\intextsep}{2pt}%
      \begin{figure}[h]
          #2
      \end{figure}
  \else

      \let\Pnotebox\Pnotebox@
      \dimen@ii\dimexpr\Psf\dimen@ii
      \setlength{\intextsep}{0pt}%
      \begin{wrapfigure}{o}[-4mm]{\dimen@ii}
          #2
      \end{wrapfigure}
      #1
      \fi
  \fi

\let\Pnotebox\Pnotebox@@
}
\makeatother



% The unformatted note text
\makeatletter
\long\def\Pnotebox@@#1#2#3{ %barcolor, backcolor, text
#3}
\makeatother

\makeatletter
\long\def\Pnotebox#1#2#3{ %barcolor, backcolor, text
#3}
\makeatother

% The formatting of the notes
\makeatletter
\newdimen\Pnoteheight
\long\def\Pnotebox@#1#2#3{ %barcolor, backcolor, text
\settoheight{\Pnoteheight}{\parbox{\dimen@ii}{#3}}%
\begin{minipage}{\dimexpr\dimen@ii+\Pcolorwidth}
  \textcolor{#1}{\rule[\dimexpr-\Pnoteheight+1mm]{\Pcolorwidth}{\dimexpr 2\Pnoteheight}}%
  \colorbox{#2}{\parbox{\dimen@ii}{#3}}
\end{minipage}%
}
\makeatother



% Commands for making the sidenote

\newwrite\tempfile
\makeatletter
\DeclareDocumentCommand \sidenote {m O{Pbardefault} O{Plightdefault}} {%
\stepcounter{localsidenotecounter}%
\hyperlink{sidenote\arabic{localsidenotecounter}}{${}^\textrm{\footnotesize\arabic{localsidenotecounter}}$}%
\immediate\write\tempfile%
{%
\noexpand\Pnotebox{#2}{#3}{%
\noexpand\footnotesize \noexpand{\noexpand\hypertarget{sidenote\arabic{localsidenotecounter}}{${}^{\arabic{localsidenotecounter}}$} \unexpanded{#1 }}%
}}}
\makeatother


%%%%HEADINGS
\usepackage{fancyhdr}
\usepackage[explicit]{titlesec}
\usepackage{lipsum}

\newsavebox\mybox
\newlength\secnumwd

\definecolor{mydarkblue}{RGB}{0,163,243}
\definecolor{mylightblue}{RGB}{191,233,251}

\titleformat{\section}
  {\normalfont\Large\sffamily\color{black}}
  {}
  {-1em}
  {%
    \savebox\mybox{\normalfont\Large\sffamily\color{black}\bfseries\thesection}%
    \settowidth\secnumwd{\usebox\mybox}%    
    \parbox[t]{\secnumwd}{{\bfseries\thesection}}\hspace{1em}%
    \parbox[t]{\dimexpr\textwidth+5em-\secnumwd-1em\relax}{#1}%
  }
  [\vskip-1.75ex\hskip-5em{\color{gray!70}\titlerule[2pt]}]
\titleformat{name=\section, numberless}
  {\normalfont\Large\sffamily\color{black}}
  {}
  {-1em}
  {#1}
  [\vskip-1.75ex\hskip-1em{\color{black}\titlerule[2pt]}]

\newcommand\FrameBoxL[1]{%
  \fcolorbox{white}{black}{\makebox[3cm][l]{\textcolor{white}{\bfseries#1}}}%
}
\newcommand\FrameBoxR[1]{%
  \fcolorbox{white}{black}{\makebox[3cm][r]{\textcolor{white}{\bfseries#1}}}%
}

\pagestyle{fancy}
%\fancyheadoffset[EL]{\dimexpr1in+\evensidemargin+\hoffset\relax}
%\fancyheadoffset[OR]{\dimexpr\paperwidth-\oddsidemargin-1in-\textwidth-\hoffset\relax}

\fancyheadoffset[EL]{45mm }
\fancyheadoffset[OR]{45mm}


\fancyhf{}
\renewcommand\headrulewidth{0pt}
\usepackage[framemethod=tikz]{mdframed}
 
 %%%%HEADINGS



%%%%EXAMPLE BOX
\definecolor{mybrown}{RGB}{128,64,0}
\mdfdefinestyle{mystyle}{%
linecolor=black,outerlinewidth=1pt,%
frametitlerule=true,frametitlefont=\sffamily\bfseries\color{white},%
frametitlerulewidth=1pt,frametitlerulecolor=black,%
frametitlebackgroundcolor=black,
backgroundcolor=mybrown!05,
innertopmargin=\topskip,
roundcorner=5pt
}
\mdtheorem[style=mystyle]{example}{Example}

\gdef\Sepline{%
  \par\noindent\makebox[\linewidth][l]{%
  \hspace*{-\mdflength{innerleftmargin}}%
   \tikz\draw[thick,dashed,gray!60] (0,0) --%
        (\textwidth+\the\mdflength{innerleftmargin}+\the\mdflength{innerrightmargin},0);
  }\par\nobreak}
%%%%EXAMPLE BOX

\definecolor{tableShade}{gray}{0.9}




\usepackage{makecell}
\newcolumntype{x}[1]{>{\centering\arraybackslash}p{#1}}
\newcolumntype{R}{>{\raggedleft\arraybackslash}X}

\usepackage{tikz}
\newcommand\diag[4]{%
  \multicolumn{1}{p{#2}|}{\hskip-\tabcolsep
  $\vcenter{\begin{tikzpicture}[baseline=0,anchor=south west,inner sep=#1]
  \path[use as bounding box] (0,0) rectangle (#2+2\tabcolsep,\baselineskip);
  \node[minimum width={#2+2\tabcolsep},minimum height=\baselineskip+\extrarowheight] (box) {};
  \draw (box.south west) -- (box.north east);
  \node[anchor=south west] at (box.south west) {#3};
  \node[anchor=north east] at (box.north east) {#4};
 \end{tikzpicture}}$\hskip-\tabcolsep}}

\definecolor{bidentitlebg}{RGB}{158,59,255}%YOUR TUBE
\newtcolorbox{ridentidad}[1][]{
  enhanced,
  frame code={
    \fill[draw=white,top color=white,bottom color=blue!60]
      ([xshift=-20pt]title.north west) --
      (title.south west) --
      (title.south east) --
      ([xshift=20pt]title.north east) -- cycle;
    \draw[blue,line width=0.4mm,rounded corners]
      (frame.north west) --
      (frame.south west) --
      ([xshift=-20pt]title.north west) --
      (title.south west) --
      (title.south east) --
      ([xshift=20pt]title.north east) --
      (frame.south east) --
      (frame.north east);
  },
  coltitle=blue!75!black,
  colback=white,
  attach boxed title to bottom center,
  boxed title style={empty},
  fonttitle=\bfseries\sffamily,
  title=\strut Visual,
  #1,
}%YOUR TUBE

\newtcolorbox{bidentidad}[1][]{%MATERIAL
  enhanced,
  skin=enhancedlast jigsaw,
  attach boxed title to top left={xshift=-4mm,yshift=-0.5mm},
  fonttitle=\bfseries\sffamily,
  colbacktitle=black,
  colframe=blue!75,
  interior style={
    top color=blue!10,
    bottom color=red!10
  },
  boxed title style={
    empty,
    arc=0pt,
    outer arc=0pt,
    boxrule=0pt
  },
  underlay boxed title={
    \fill[blue!75!white] 
      (title.north west) -- 
      (title.north east) -- 
      +(\tcboxedtitleheight-1mm,-\tcboxedtitleheight+1mm) -- 
      ([xshift=4mm,yshift=0.5mm]frame.north east) -- 
      +(0mm,-1mm) -- 
      (title.south west) -- cycle;
    \fill[blue!45!white!50!black] 
      ([yshift=-0.5mm]frame.north west) -- 
      +(-0.4,0) -- 
      +(0,-0.3) -- cycle;
    \fill[blue!45!white!50!black] 
      ([yshift=-0.5mm]frame.north east) -- 
      +(0,-0.3) -- 
      +(0.4,0) -- cycle; 
  },
  title={Material},
  #1
}%MATERIAL
%GRAPH PAPER
\usepackage{graphpap}
\makeatletter
\newcommand\graphpapern[1][10]{\leavevmode\@gridn{#1}}
\def\@gridn#1(#2,#3)#4{\@grid@n{#1}{#2}{#3}(}
\def\@grid@n#1#2#3(#4,#5){%
\@tempcnta=#4\relax
\divide\@tempcnta#1\relax
\advance\@tempcnta1\relax
{\thinlines\@nonumvgrid(#2,#3){#1}{\@tempcnta}{#5}
\@tempcnta#4\relax
\divide\@tempcnta5\relax
\divide\@tempcnta#1\relax
\advance\@tempcnta1\relax
\@tempcntb5\relax
\multiply\@tempcntb#1\relax
\thicklines\@nonumvgrid(#2,#3){\@tempcntb}{\@tempcnta}{#5}
\@tempcnta#5\relax
\divide\@tempcnta #1\relax
\advance\@tempcnta1\relax
\thinlines\@nonumhgrid(#2,#3){#1}{\@tempcnta}{#4}
\@tempcnta#5\relax
\divide\@tempcnta5\relax
\divide\@tempcnta#1\relax
\advance\@tempcnta1\relax
\thicklines\@nonumhgrid(#2,#3){\@tempcntb}{\@tempcnta}{#4}}%
\ignorespaces}
\makeatother
% SIDE TABLE
 \usepackage{rotating,booktabs,dcolumn}
\usepackage[subpreambles=true]{standalone}
\usepackage{import}
\usepackage{tabu}
\usepackage{tabularx} % in the preamble
\usepackage{pgfplots}
\usepackage{cancel} 
\taburulecolor{black}

%PAGE NUMBER
\fancyhead[OR]{\raisebox{1mm}[0pt][0pt]{  \thepage } }
\fancyhead[EL]{\FrameBoxR{\thepage}}
%%%%%%%%%%%%%%%%%%%%%%%%%%%%%%%%%%%%%%%%%%%%%%%%%%%
%%%%%%%%%%%%%%%%%%%%%%%%%%%%%%%%%%%%%%%%%%%%%%%%%%%
%%%%%%%%%%%%%%%%%%%%%%%%%%%%%%%%%%%%%%%%%%%%%%%%%%%
%%%%%%%%%%%%%%%%%%%%%%%%%%%%%%%%%%%%%%%%%%%%%%%%%%%
%%%%%%%%%%%%%%%%%%%%%%%%%%%%%%%%%%%%%%%%%%%%%%%%%%%
\usepackage{makeidx} 
\makeindex
\renewcommand*\contentsname{Summary}
\usepackage{tocloft}
\renewcommand\cftchapaftersnum{.}
\renewcommand\cftsecaftersnum{.}
\renewcommand\thechapter{\Roman{chapter}}
\renewcommand\thesection{\arabic{section}}
\renewcommand\thesubsection{(\alph{subsection})}

%\usepackage{textcomp}
%\usepackage{xcolor}
%\usepackage{soul}
%\usepackage{capt-of}
%\pgfplotsset{compat=newest}
\usepackage{pbox}
%\usepackage[svgnames]{xcolor} % Required to specify font color

\usepackage{adjustbox}


\usepackage{pdfpages}


\newcommand\myfigure[5]{%
  \ifdim#2>.8\linewidth
    {%
      \centering
      \includegraphics[width=#3]{#4}%
%      \captionof{figure}{#5}%
    }%
  \else
  \begin{wrapfigure}{#1}{#2}
    \includegraphics[width=#3]{#4}
%    \caption{#5}
  \end{wrapfigure}
  \fi
}



%HIGHLIGHT IN MATH
\newcommand{\hlmath}[1]{%
  \colorbox{black!20}{$\displaystyle#1$}}
  %HIGHLIGHT IN MATH


%\usepackage[backend=bibtex,style=chem-acs,citestyle=numeric-comp, sorting=none, defernums=true, hyperref, backref, refsegment=chapter]{biblatex}
%\usepackage[style=chem-acs, citestyle=numeric-comp,backend=bibtex, refsegment=chapter, defernumbers=true]{biblatex}
 \usepackage[style=chem-acs, backend=bibtex]{biblatex}

\addbibresource{./biblio/bibliography.bib}

%%%% CHECKBOX
\usepackage{amssymb,tikz,pdftexcmds,xparse}
\tikzset{box/.style={
    minimum size=0.225cm,
    inner sep=0pt,
    draw,
  },  
  insert mark/.style={
    append after command={%
         node[inner sep=0pt,#1]
           at (\tikzlastnode.center){$\checkmark$}
     }     
  },
  insert bad mark/.style={
    append after command={%
         [shorten <=\pgflinewidth,shorten >=\pgflinewidth]
         (\tikzlastnode.north west)edge[#1](\tikzlastnode.south east)
         (\tikzlastnode.south west)edge[#1](\tikzlastnode.north east)
     }     
  },
}
\makeatletter
\NewDocumentCommand{\tikzcheckmark}{O{} m}{%
  \ifnum\pdf@strcmp{#2}{mark}=\z@%
    \tikz[baseline=-0.5ex]\node[box,insert mark={#1},#1]{};%
  \fi%
  \ifnum\pdf@strcmp{#2}{bad mark}=\z@%
    \tikz[baseline=-0.5ex]\node[box,insert bad mark={#1},#1]{};%
  \fi%
  \ifnum\pdf@strcmp{#2}{no mark}=\z@%
    \tikz[baseline=-0.5ex]\node[box,#1]{};%
  \fi%
}
\makeatother
%%%%CHECKBOX

 
%%%%% NEW TABLES OF CONTENT
\definecolor{doc}{RGB}{0,60,110}
%\contentsmargin{0cm}
\usepackage{titletoc}
\titlecontents{part}[0pc]
{\addvspace{30pt}%
\begin{tikzpicture}[remember picture, overlay]%
\draw[fill=blue!75,draw=blue!75] (+1.9,-.1) rectangle (0,.5);%
\pgftext[left,x=0.5cm,y=0.2cm]{\color{white}\Large\sc\bfseries part\ \hspace*{4cm} \thecontentslabel};%
\end{tikzpicture}\color{blue!75}\large\sc\bfseries}%
{}
{}
{\;\titlerule\;\large\sc\bfseries Page \thecontentspage
}%
\titlecontents{section}[2.4pc]
{\addvspace{1pt}}
{\contentslabel[\thecontentslabel]{2.4pc}}
{}
{\hfill\small \thecontentspage}
[]




\titlecontents{section}[2.4pc]
{\addvspace{1pt}}
{\contentslabel[\thecontentslabel]{2.4pc}}
{}
{\hfill\small \thecontentspage}
[]
\titlecontents*{subsection}[4pc]
{\addvspace{-1pt}\small}
{}
{}
{\ --- \small\thecontentspage}
[ \textbullet\ ][]

\makeatletter
\renewcommand{\tableofcontents}{%
\chapter*{%
\vspace*{-20\p@}%
\begin{tikzpicture}[remember picture, overlay]%
\pgftext[right,x=-2cm,y=0.2cm]{\color{doc!30}\Huge\sc\bfseries \contentsname};%
\draw[fill=doc!30,draw=doc!30] (1,-.75) rectangle (20,1);%
\clip (13,-.75) rectangle (20,1);
\pgftext[right,x=-5cm,y=0.2cm]{\color{red}\Huge\sc\bfseries \contentsname};%
\end{tikzpicture}}%
\@starttoc{toc}}
\makeatother


\usepackage{everypage}%%% SIDEBAR ON PAGES
\usepackage{xcolor}
\usepackage{scrlayer}
\DeclareNewLayer[
  background,
  outermargin,
  addhoffset=\marginparsep,
  contents={\color{gray!50}\rule{\layerwidth}{\layerheight}}
]{sidebar.bg}
\makeatletter
\DeclareNewLayer[
  clone=sidebar.bg,
  addhoffset=1em,
  addwidth=-1em,
  contents={\parbox[b][\layerheight]{\layerwidth}{\raggedright\@sidebar}}
]{sidebar.txt}
\newcommand\@sidebar{}
\newcommand\sidebar[1]{\renewcommand\@sidebar{#1}}
\makeatother
\DeclareNewPageStyleByLayers{sidebar}{sidebar.bg,sidebar.txt}


%%%%% COLOR TABLE
\usepackage{environ}
\usepackage{tikz}
\usepackage{longtable}
\newcommand*\rot{\rotatebox{90}}
%\usepackage{chemfig}
\def\Z{  0 & 0 & 12 & Working with leaflets & Claudio Fiandrino\\}
\def\ZZ{\Z\Z\Z\Z}\def\ZZZ{\ZZ\ZZ\ZZ\ZZ}\def\ZZZZ{\ZZZ\ZZZ\ZZZ\ZZZ}

% some stuff common to short and long tables
\definecolor{tablecolor}{named}{ForestGreen}
\newcommand{\header}[1]{%
  \multicolumn{1}{c}{%
    \cellcolor{tablecolor}\color{white}\bfseries#1}}
\newcommand{\rndtablesetup}{%
  \addtolength{\extrarowheight}{1ex}%
  \rowcolors{2}{tablecolor!20}{tablecolor!40}%
  \sffamily}

% styles for the clipping/internal nodes
\tikzset{
  myclip/.style={rounded corners=1ex},
  mynode/.style={anchor=south west,inner sep=0pt,outer sep=0pt},
  mybord/.style={myclip,very thick,tablecolor},
}

% the short table (provided)
\newsavebox{\tablebox}
\NewEnviron{rndtable}[1]{%
  \rndtablesetup%
  \savebox{\tablebox}{%
    \begin{tabular}{#1}%
      \BODY%
    \end{tabular}}%
  \begin{tikzpicture}
    \clip[myclip] (0,-\dp\tablebox) -- (\wd\tablebox,-\dp\tablebox) --
      (\wd\tablebox,\ht\tablebox) -- (0,\ht\tablebox) -- cycle;
    \node[mynode] at (0,-\dp\tablebox) {\usebox{\tablebox}};
    \draw[mybord] (0,-\dp\tablebox) -- (\wd\tablebox,-\dp\tablebox) --
      (\wd\tablebox,\ht\tablebox) -- (0,\ht\tablebox) -- cycle;
  \end{tikzpicture}}

% round top
\newcommand{\roundtop}[1]{%
  \begin{tikzpicture}
    \clip[myclip] (0,0) |- (\wd#1,\ht#1+\dp#1) -- (\wd#1,0);
    \draw[mybord] (0,0) |- (\wd#1,\ht#1+\dp#1) -- (\wd#1,0);
    \node[mynode] at (0,0) {\box#1};
  \end{tikzpicture}}

% round bottom
\newcommand{\roundbot}[1]{%
  \begin{tikzpicture}
    \clip[myclip] (0,0) |- (\wd#1,-\ht#1+\dp#1) -- (\wd#1,0);
    \draw[mybord] (0,0) |- (\wd#1,-\ht#1+\dp#1) -- (\wd#1,0);
    \node[mynode] at (0,0) {\box#1};
  \end{tikzpicture}}

% frame midsections
\newcommand{\framemid}[1]{%
  \begin{tikzpicture}
    \draw[mybord] (0,0) -- (0,\ht#1+\dp#1) (\wd#1,0) -- (\wd#1,\ht#1+\dp#1);
    \node[mynode] at (0,0) {\box#1};
  \end{tikzpicture}}




\usepackage[font=small]{caption}% no need in this mwe
\usepackage{float}% no need in this mwe
\usepackage{floatflt}% no need in this mwe
\usepackage{tikz}
%\usepackage{tikzorbital}
\usetikzlibrary{fadings,patterns,backgrounds,fit,arrows}
\pgfdeclarelayer{backbackground}
\pgfsetlayers{backbackground,background,main}
%\oddsidemargin=-10.4mm  %
%\evensidemargin=-20.4mm %
%\topmargin=-24mm        %
%\textwidth=190mm        %
%\textheight=275mm       %
%\headheight = 5mm      %
%\headsep = 4mm
%\footskip = 8mm


%% CIRCLE LEWIS
\newcommand*\circleatom[1]{\tikz\node[circle, fill=green!30]{\printatom{#1}};}
%% CIRCLE LEWIS

%%% ORBITAL DIAGRAMS
\usepackage{amssymb} % for harpoons
\newcommand{\electron}[2]{{%
        \newcommand*\up{\fbox{$\mathord\upharpoonleft\phantom{\downharpoonright}$}}%
        \newcommand*\dwn{\fbox{$\mathord\downharpoonleft\phantom{\upharpoonright}$}}%
        \newcommand*\updwn{\fbox{$\upharpoonleft\downharpoonright$}}%
        \newcommand*\emp{\fbox{$\phantom{\downharpoonright}\phantom{\downharpoonright}$}}%
        \setlength\tabcolsep{0pt}% remove extra horizontal space from tabular
        \begin{tabular}{c}#2\\[2pt]#1\end{tabular}}}%
%%% ORBITAL DIAGRAMS

\newcommand*{\arraycolor}[1]{\protect\leavevmode\color{#1}}
\newcolumntype{A}{>{\columncolor{blue!50!white}    }c}
%\newcolumntype{B}{>{\columncolor{LightGoldenrod}}c}
%\newcolumntype{C}{>{\columncolor{FireBrick!50}}c}
%\newcolumntype{D}{>{\columncolor{Gray!42}}c}
\newcolumntype{D}[1]{>{\columncolor{Gray!42} \centering\let\newline\\\arraybackslash\hspace{0pt}}m{#1}}
\newcolumntype{E}{>{\columncolor{white}}c}
\newcolumntype{F}[1]{>{\columncolor{blue!50!white} \centering\let\newline\\\arraybackslash\hspace{0pt}}m{#1}}
\newcolumntype{G}{>{\columncolor{Gray!42}}c}
\newcolumntype{B}[1]{>{\columncolor{FireBrick!50} \centering\let\newline\\\arraybackslash\hspace{0pt}}m{#1}}
\newcolumntype{C}[1]{>{\columncolor{LightGoldenrod} \centering\let\newline\\\arraybackslash\hspace{0pt}}m{#1}}
\newcolumntype{E}[1]{>{\columncolor{Gray!22} \centering\let\newline\\\arraybackslash\hspace{0pt}}m{#1}}
\newcolumntype{M}[1]{>{\centering\arraybackslash}m{#1}}

\newcommand\tikzmarkA[2]{%
\tikz[remember picture,baseline] \node[inner sep=2pt,outer sep=0] (#1){#2};%
}

\newcommand\link[2]{%
\begin{tikzpicture}[remember picture, overlay, >=stealth, shift={(0,0)}]
  \draw[->] (#1) to (#2);
\end{tikzpicture}%
}



\usetikzlibrary{fadings,patterns,fit}

% use stamp.png to make a fading (a mask)
\begin{tikzfadingfrompicture}[name=stamp]
  \node [fill=transparent!0,inner sep=0]
  {\includegraphics[width=50pt,height=50pt]{./stamp.png}};
  \begin{scope}[yshift=-90pt,transform canvas={scale=.2}]
    \node[color=gray!50!black,font=\ttfamily] {\today};
  \end{scope}
\end{tikzfadingfrompicture}

\pgfdeclarehorizontalshading[mycolor]{stampshading}
  {4cm}{
    color(0cm)=(mycolor);
    color(1cm)=(mycolor);
    color(1.35cm)=(mycolor);
    color(1.40cm)=(white!80!mycolor);
    color(1.45cm)=(mycolor);
    color(1.85cm)=(mycolor);
    color(2.00cm)=(white!80!mycolor);
    color(2.05cm)=(mycolor);
    color(2.25cm)=(mycolor);
    color(2.30cm)=(white!80!mycolor);
    color(2.35cm)=(mycolor);
    color(2.45cm)=(mycolor);
    color(2.50cm)=(white!80!mycolor);
    color(2.55cm)=(mycolor);
    color(4cm)=(mycolor)
  }
  
  \usepackage{listings}
  \usepackage{caption}
\DeclareCaptionFont{white}{\color{white}}
\DeclareCaptionFormat{listing}{\colorbox{lightgray}{\parbox{\textwidth}{#1#2#3}}}
\captionsetup[lstlisting]{format=listing,labelfont=white,textfont=white}
\usepackage{longtable}




%%%%%EMTPY PERIDOC TABLE
%%%%%%%%%%%%%%%%%%%
%%%%%%%%%%%%%%%%%%%

\newcommand{\CommonElementTextFormat}[4]
{
  \begin{minipage}{2.2cm}
    \centering
      {\textbf{#1} \hfill #2}%
      \linebreak \linebreak
      {\textbf{#3}}%
      \linebreak \linebreak
      {{#4}}
  \end{minipage}
}

\newcommand{\NaturalElementTextFormat}[4]
{
  \CommonElementTextFormat{#1}{#2}{\Huge {#3}}{#4}
}

\newcommand{\OutlineText}[1]
{
\ifpdf
  % Couldn't find a nicer way of doing an outline font with TikZ
  % other than using pdfliteral 1 Tr
  %
  \pdfliteral direct {0.5 w 1 Tr}{#1}%
  \pdfliteral direct {1 w 0 Tr}%
\else
  % pstricks can do this with \pscharpath from pstricks
  %
  \pscharpath[shadow=false,
    fillstyle=solid,
    fillcolor=white,
    linestyle=solid,
    linecolor=black,
    linewidth=.2pt]{#1} 
\fi
}

\newcommand{\SyntheticElementTextFormat}[4]
{
\ifpdf
  \CommonElementTextFormat{#1}{#2}{\OutlineText{\Huge #3}}{#4}
\else
  % pstricks approach results in slightly larger box
  % that doesn't break, so fudge here
  \CommonElementTextFormat{#1}{#2}{\OutlineText{\Large #3}}{#4}
\fi
}


\newcommand{\PeriodicTableFull}{
\begin{tikzpicture}[font=\sffamily, scale=0.30, transform shape]

%% Fill Color Styles
  \tikzstyle{ElementFill} = [fill=yellow!15]
  \tikzstyle{AlkaliMetalFill} = [fill=blue!55]
  \tikzstyle{AlkalineEarthMetalFill} = [fill=blue!40]
  \tikzstyle{MetalFill} = [fill=blue!25]
  \tikzstyle{MetalloidFill} = [fill=orange!25]
  \tikzstyle{NonmetalFill} = [fill=green!5]
  \tikzstyle{ChalcogenFill} = [fill=red!25]
  \tikzstyle{HalogenFill} = [fill=green!40]
  \tikzstyle{NobleGasFill} = [fill=green!55]
  \tikzstyle{LanthanideActinideFill} = [fill=purple!25]

%% Element Styles
  \tikzstyle{Element} = [draw=black, ElementFill,
    minimum width=2.75cm, minimum height=2.75cm, node distance=2.75cm]
  \tikzstyle{AlkaliMetal} = [Element, AlkaliMetalFill]
  \tikzstyle{AlkalineEarthMetal} = [Element, AlkalineEarthMetalFill]
  \tikzstyle{Metal} = [Element, MetalFill]
  \tikzstyle{Metalloid} = [Element, MetalloidFill]
  \tikzstyle{Nonmetal} = [Element, NonmetalFill]
  \tikzstyle{Chalcogen} = [Element, ChalcogenFill]
  \tikzstyle{Halogen} = [Element, HalogenFill]
  \tikzstyle{NobleGas} = [Element, NobleGasFill]
  \tikzstyle{LanthanideActinide} = [Element, LanthanideActinideFill]
  \tikzstyle{PeriodLabel} = [font={\sffamily\LARGE}, node distance=2.0cm]
  \tikzstyle{GroupLabel} = [font={\sffamily\LARGE}, minimum width=2.75cm, node distance=2.0cm]
  \tikzstyle{TitleLabel} = [font={\sffamily\Huge\bfseries}]

%% Group 1 - IA
  \node[name=H, Element] {\NaturalElementTextFormat{1}{1.0079}{H}{Hydrogen}};
  \node[name=Li, below of=H, AlkaliMetal] {\NaturalElementTextFormat{3}{6.941}{Li}{Lithium}};
  \node[name=Na, below of=Li, AlkaliMetal] {\NaturalElementTextFormat{11}{22.990}{Na}{Sodium}};
  \node[name=K, below of=Na, AlkaliMetal] {\NaturalElementTextFormat{19}{39.098}{K}{Potassium}};
  \node[name=Rb, below of=K, AlkaliMetal] {\NaturalElementTextFormat{37}{85.468}{Rb}{Rubidium}};
  \node[name=Cs, below of=Rb, AlkaliMetal] {\NaturalElementTextFormat{55}{132.91}{Cs}{Caesium}};
  \node[name=Fr, below of=Cs, AlkaliMetal] {\NaturalElementTextFormat{87}{223}{Fr}{Francium}};

%% Group 2 - IIA
  \node[name=Be, right of=Li, AlkalineEarthMetal] {\NaturalElementTextFormat{4}{9.0122}{Be}{Beryllium}};
  \node[name=Mg, below of=Be, AlkalineEarthMetal] {\NaturalElementTextFormat{12}{24.305}{Mg}{Magnesium}};
  \node[name=Ca, below of=Mg, AlkalineEarthMetal] {\NaturalElementTextFormat{20}{40.078}{Ca}{Calcium}};
  \node[name=Sr, below of=Ca, AlkalineEarthMetal] {\NaturalElementTextFormat{38}{87.62}{Sr}{Strontium}};
  \node[name=Ba, below of=Sr, AlkalineEarthMetal] {\NaturalElementTextFormat{56}{137.33}{Ba}{Barium}};
  \node[name=Ra, below of=Ba, AlkalineEarthMetal] {\NaturalElementTextFormat{88}{226}{Ra}{Radium}};

%% Group 3 - IIIB
  \node[name=Sc, right of=Ca, Metal] {\NaturalElementTextFormat{21}{44.956}{Sc}{Scandium}};
  \node[name=Y, below of=Sc, Metal] {\NaturalElementTextFormat{39}{88.906}{Y}{Yttrium}};
  \node[name=LaLu, below of=Y, LanthanideActinide] {\NaturalElementTextFormat{57-71}{}{La-Lu}{Lanthanide}};
  \node[name=AcLr, below of=LaLu, LanthanideActinide] {\NaturalElementTextFormat{89-103}{}{Ac-Lr}{Actinide}};

%% Group 4 - IVB
  \node[name=Ti, right of=Sc, Metal] {\NaturalElementTextFormat{22}{47.867}{Ti}{Titanium}};
  \node[name=Zr, below of=Ti, Metal] {\NaturalElementTextFormat{40}{91.224}{Zr}{Zirconium}};
  \node[name=Hf, below of=Zr, Metal] {\NaturalElementTextFormat{72}{178.49}{Hf}{Halfnium}};
  \node[name=Rf, below of=Hf, Metal] {\SyntheticElementTextFormat{104}{261}{Rf}{Rutherfordium}};

%% Group 5 - VB
  \node[name=V, right of=Ti, Metal] {\NaturalElementTextFormat{23}{50.942}{V}{Vanadium}};
  \node[name=Nb, below of=V, Metal] {\NaturalElementTextFormat{41}{92.906}{Nb}{Niobium}};
  \node[name=Ta, below of=Nb, Metal] {\NaturalElementTextFormat{73}{180.95}{Ta}{Tantalum}};
  \node[name=Db, below of=Ta, Metal] {\SyntheticElementTextFormat{105}{262}{Db}{Dubnium}};

%% Group 6 - VIB
  \node[name=Cr, right of=V, Metal] {\NaturalElementTextFormat{24}{51.996}{Cr}{Chromium}};
  \node[name=Mo, below of=Cr, Metal] {\NaturalElementTextFormat{42}{95.94}{Mo}{Molybdenum}};
  \node[name=W, below of=Mo, Metal] {\NaturalElementTextFormat{74}{183.84}{W}{Tungsten}};
  \node[name=Sg, below of=W, Metal] {\SyntheticElementTextFormat{106}{266}{Sg}{Seaborgium}};

%% Group 7 - VIIB
  \node[name=Mn, right of=Cr, Metal] {\NaturalElementTextFormat{25}{54.938}{Mn}{Manganese}};
  \node[name=Tc, below of=Mn, Metal] {\NaturalElementTextFormat{43}{96}{Tc}{Technetium}};
  \node[name=Re, below of=Tc, Metal] {\NaturalElementTextFormat{75}{186.21}{Re}{Rhenium}};
  \node[name=Bh, below of=Re, Metal] {\SyntheticElementTextFormat{107}{264}{Bh}{Bohrium}};

%% Group 8 - VIIIB
  \node[name=Fe, right of=Mn, Metal] {\NaturalElementTextFormat{26}{55.845}{Fe}{Iron}};
  \node[name=Ru, below of=Fe, Metal] {\NaturalElementTextFormat{44}{101.07}{Ru}{Ruthenium}};
  \node[name=Os, below of=Ru, Metal] {\NaturalElementTextFormat{76}{190.23}{Os}{Osmium}};
  \node[name=Hs, below of=Os, Metal] {\SyntheticElementTextFormat{108}{277}{Hs}{Hassium}};

%% Group 9 - VIIIB
  \node[name=Co, right of=Fe, Metal] {\NaturalElementTextFormat{27}{58.933}{Co}{Cobalt}};
  \node[name=Rh, below of=Co, Metal] {\NaturalElementTextFormat{45}{102.91}{Rh}{Rhodium}};
  \node[name=Ir, below of=Rh, Metal] {\NaturalElementTextFormat{77}{192.22}{Ir}{Iridium}};
  \node[name=Mt, below of=Ir, Metal] {\SyntheticElementTextFormat{109}{268}{Mt}{Meitnerium}};

%% Group 10 - VIIIB
  \node[name=Ni, right of=Co, Metal] {\NaturalElementTextFormat{28}{58.693}{Ni}{Nickel}};
  \node[name=Pd, below of=Ni, Metal] {\NaturalElementTextFormat{46}{106.42}{Pd}{Palladium}};
  \node[name=Pt, below of=Pd, Metal] {\NaturalElementTextFormat{78}{195.08}{Pt}{Platinum}};
  \node[name=Ds, below of=Pt, Metal] {\SyntheticElementTextFormat{110}{281}{Ds}{Darmstadtium}};

%% Group 11 - IB
  \node[name=Cu, right of=Ni, Metal] {\NaturalElementTextFormat{29}{63.546}{Cu}{Copper}};
  \node[name=Ag, below of=Cu, Metal] {\NaturalElementTextFormat{47}{107.87}{Ag}{Silver}};
  \node[name=Au, below of=Ag, Metal] {\NaturalElementTextFormat{79}{196.97}{Au}{Gold}};
  \node[name=Rg, below of=Au, Metal] {\SyntheticElementTextFormat{111}{280}{Rg}{Roentgenium}};

%% Group 12 - IIB
  \node[name=Zn, right of=Cu, Metal] {\NaturalElementTextFormat{30}{65.39}{Zn}{Zinc}};
  \node[name=Cd, below of=Zn, Metal] {\NaturalElementTextFormat{48}{112.41}{Cd}{Cadmium}};
  \node[name=Hg, below of=Cd, Metal] {\NaturalElementTextFormat{80}{200.59}{Hg}{Mercury}};
  \node[name=Uub, below of=Hg, Metal] {\SyntheticElementTextFormat{112}{285}{Uub}{Ununbium}};

%% Group 13 - IIIA
  \node[name=Ga, right of=Zn, Metal] {\NaturalElementTextFormat{31}{69.723}{Ga}{Gallium}};
  \node[name=Al, above of=Ga, Metal] {\NaturalElementTextFormat{13}{26.982}{Al}{Aluminium}};
  \node[name=B, above of=Al, Metalloid] {\NaturalElementTextFormat{5}{10.811}{B}{Boron}};
  \node[name=In, below of=Ga, Metal] {\NaturalElementTextFormat{49}{114.82}{In}{Indium}};
  \node[name=Tl, below of=In, Metal] {\NaturalElementTextFormat{81}{204.38}{Tl}{Thallium}};
  \node[name=Uut, below of=Tl, Metal] {\SyntheticElementTextFormat{113}{284}{Uut}{Ununtrium}};

%% Group 14 - IVA
  \node[name=C, right of=B, Nonmetal] {\NaturalElementTextFormat{6}{12.011}{C}{Carbon}};
  \node[name=Si, below of=C, Metalloid] {\NaturalElementTextFormat{14}{28.086}{Si}{Silicon}};
  \node[name=Ge, below of=Si, Metalloid] {\NaturalElementTextFormat{32}{72.64}{Ge}{Germanium}};
  \node[name=Sn, below of=Ge, Metal] {\NaturalElementTextFormat{50}{118.71}{Sn}{Tin}};
  \node[name=Pb, below of=Sn, Metal] {\NaturalElementTextFormat{82}{207.2}{Pb}{Lead}};
  \node[name=Uuq, below of=Pb, Metal] {\SyntheticElementTextFormat{114}{289}{Uuq}{Ununquadium}};

%% Group 15 - VA
  \node[name=N, right of=C, Nonmetal] {\NaturalElementTextFormat{7}{14.007}{N}{Nitrogen}};
  \node[name=P, below of=N, Nonmetal] {\NaturalElementTextFormat{15}{30.974}{P}{Phosphorus}};
  \node[name=As, below of=P, Metalloid] {\NaturalElementTextFormat{33}{74.922}{As}{Arsenic}};
  \node[name=Sb, below of=As, Metalloid] {\NaturalElementTextFormat{51}{121.76}{Sb}{Antimony}};
  \node[name=Bi, below of=Sb, Metal] {\NaturalElementTextFormat{83}{208.98}{Bi}{Bismuth}};
  \node[name=Uup, below of=Bi, Metal] {\SyntheticElementTextFormat{115}{288}{Uup}{Ununpentium}};

%% Group 16 - VIA
  \node[name=O, right of=N, Chalcogen] {\NaturalElementTextFormat{8}{15.999}{O}{Oxygen}};
  \node[name=S, below of=O, Chalcogen] {\NaturalElementTextFormat{16}{32.065}{S}{Sulphur}};
  \node[name=Se, below of=S, Chalcogen] {\NaturalElementTextFormat{34}{78.96}{Se}{Selenium}};
  \node[name=Te, below of=Se, Metalloid] {\NaturalElementTextFormat{52}{127.6}{Te}{Tellurium}};
  \node[name=Po, below of=Te, Metalloid] {\NaturalElementTextFormat{84}{209}{Po}{Polonium}};
  \node[name=Uuh, below of=Po, Metal] {\SyntheticElementTextFormat{116}{293}{Uuh}{Ununhexium}};
%  %% Group 16 - VIA
%  \node[name=O, right of=N, Nonmetal] {\NaturalElementTextFormat{8}{15.999}{O}{Oxygen}};
%  \node[name=S, below of=O, Nonmetal] {\NaturalElementTextFormat{16}{32.065}{S}{Sulphur}};
%  \node[name=Se, below of=S, Nonmetal] {\NaturalElementTextFormat{34}{78.96}{Se}{Selenium}};
%  \node[name=Te, below of=Se, Metalloid] {\NaturalElementTextFormat{52}{127.6}{Te}{Tellurium}};
%  \node[name=Po, below of=Te, Metalloid] {\NaturalElementTextFormat{84}{209}{Po}{Polonium}};
%  \node[name=Uuh, below of=Po, Metal] {\SyntheticElementTextFormat{116}{293}{Uuh}{Ununhexium}};
%
%% Group 17 - VIIA
  \node[name=F, right of=O, Halogen] {\NaturalElementTextFormat{9}{18.998}{F}{Flourine}};
  \node[name=Cl, below of=F, Halogen] {\NaturalElementTextFormat{17}{35.453}{Cl}{Chlorine}};
  \node[name=Br, below of=Cl, Halogen] {\NaturalElementTextFormat{35}{79.904}{Br}{Bromine}};
  \node[name=I, below of=Br, Halogen] {\NaturalElementTextFormat{53}{126.9}{I}{Iodine}};
  \node[name=At, below of=I, Metalloid] {\NaturalElementTextFormat{85}{210}{At}{Astatine}};
  \node[name=Uus, below of=At, Element] {\SyntheticElementTextFormat{117}{292}{Uus}{Ununseptium}}; 

%% Group 18 - VIIIA
  \node[name=Ne, right of=F, NobleGas] {\NaturalElementTextFormat{10}{20.180}{Ne}{Neon}};
  \node[name=He, above of=Ne, NobleGas] {\NaturalElementTextFormat{2}{4.0025}{He}{Helium}};
  \node[name=Ar, below of=Ne, NobleGas] {\NaturalElementTextFormat{18}{39.948}{Ar}{Argon}};
  \node[name=Kr, below of=Ar, NobleGas] {\NaturalElementTextFormat{36}{83.8}{Kr}{Krypton}};
  \node[name=Xe, below of=Kr, NobleGas] {\NaturalElementTextFormat{54}{131.29}{Xe}{Xenon}};
  \node[name=Rn, below of=Xe, NobleGas] {\NaturalElementTextFormat{86}{222}{Rn}{Radon}};
  \node[name=Uuo, below of=Rn, Nonmetal] {\SyntheticElementTextFormat{118}{294}{Uuo}{Ununoctium}}; 

%% Period
  \node[name=Period1, left of=H, PeriodLabel] {1};
  \node[name=Period2, left of=Li, PeriodLabel] {2};
  \node[name=Period3, left of=Na, PeriodLabel] {3}; 
  \node[name=Period4, left of=K, PeriodLabel] {4}; 
  \node[name=Period5, left of=Rb, PeriodLabel] {5};
  \node[name=Period6, left of=Cs, PeriodLabel] {6};
  \node[name=Period7, left of=Fr, PeriodLabel] {7};

%% Group
  \node[name=Group1, above of=H, GroupLabel] {1 \hfill IA};
  \node[name=Group2, above of=Be, GroupLabel] {2 \hfill IIA};
  \node[name=Group3, above of=Sc, GroupLabel] {3 \hfill IIIB};
  \node[name=Group4, above of=Ti, GroupLabel] {4 \hfill IVB};
  \node[name=Group5, above of=V, GroupLabel] {5 \hfill VB};
  \node[name=Group6, above of=Cr, GroupLabel] {6 \hfill VIB};
  \node[name=Group7, above of=Mn, GroupLabel] {7 \hfill VIIB};
  \node[name=Group8, above of=Fe, GroupLabel] {8 \hfill VIIIB};
  \node[name=Group9, above of=Co, GroupLabel] {9 \hfill VIIIB};
  \node[name=Group10, above of=Ni, GroupLabel] {10 \hfill VIIIB};
  \node[name=Group11, above of=Cu, GroupLabel] {11 \hfill IB};
  \node[name=Group12, above of=Zn, GroupLabel] {12 \hfill IIB};
  \node[name=Group13, above of=B, GroupLabel] {13 \hfill IIIA};
  \node[name=Group14, above of=C, GroupLabel] {14 \hfill IVA};
  \node[name=Group15, above of=N, GroupLabel] {15 \hfill VA};
  \node[name=Group16, above of=O, GroupLabel] {16 \hfill VIA};
  \node[name=Group17, above of=F, GroupLabel] {17 \hfill VIIA};
  \node[name=Group18, above of=He, GroupLabel] {18 \hfill VIIIA};

%% Lanthanide
  \node[name=La, below of=Rf, LanthanideActinide, yshift=-1cm] {\NaturalElementTextFormat{57}{138.91}{La}{Lanthanum}};
  \node[name=Ce, right of=La, LanthanideActinide] {\NaturalElementTextFormat{58}{140.12}{Ce}{Cerium}};
  \node[name=Pr, right of=Ce, LanthanideActinide] {\NaturalElementTextFormat{59}{140.91}{Pr}{Praseodymium}};
  \node[name=Nd, right of=Pr, LanthanideActinide] {\NaturalElementTextFormat{60}{144.24}{Nd}{Neodymium}};
  \node[name=Pm, right of=Nd, LanthanideActinide] {\NaturalElementTextFormat{61}{145}{Pm}{Promethium}};
  \node[name=Sm, right of=Pm, LanthanideActinide] {\NaturalElementTextFormat{62}{150.36}{Sm}{Samarium}};
  \node[name=Eu, right of=Sm, LanthanideActinide] {\NaturalElementTextFormat{63}{151.96}{Eu}{Europium}};
  \node[name=Gd, right of=Eu, LanthanideActinide] {\NaturalElementTextFormat{64}{157.25}{Gd}{Gadolinium}};
  \node[name=Tb, right of=Gd, LanthanideActinide] {\NaturalElementTextFormat{65}{158.93}{Tb}{Terbium}};
  \node[name=Dy, right of=Tb, LanthanideActinide] {\NaturalElementTextFormat{66}{162.50}{Dy}{Dysprosium}};
  \node[name=Ho, right of=Dy, LanthanideActinide] {\NaturalElementTextFormat{67}{164.93}{Ho}{Holmium}};
  \node[name=Er, right of=Ho, LanthanideActinide] {\NaturalElementTextFormat{68}{167.26}{Er}{Erbium}};
  \node[name=Tm, right of=Er, LanthanideActinide] {\NaturalElementTextFormat{69}{168.93}{Tm}{Thulium}};
  \node[name=Yb, right of=Tm, LanthanideActinide] {\NaturalElementTextFormat{70}{173.04}{Yb}{Ytterbium}};
  \node[name=Lu, right of=Yb, LanthanideActinide] {\NaturalElementTextFormat{71}{174.97}{Lu}{Lutetium}};

%% Actinide
  \node[name=Ac, below of=La, LanthanideActinide, yshift=-1cm] {\NaturalElementTextFormat{89}{227}{Ac}{Actinium}};
  \node[name=Th, right of=Ac, LanthanideActinide] {\NaturalElementTextFormat{90}{232.04}{Th}{Thorium}};
  \node[name=Pa, right of=Th, LanthanideActinide] {\NaturalElementTextFormat{91}{231.04}{Pa}{Protactinium}};
  \node[name=U, right of=Pa, LanthanideActinide] {\NaturalElementTextFormat{92}{238.03}{U}{Uranium}};
  \node[name=Np, right of=U, LanthanideActinide] {\SyntheticElementTextFormat{93}{237}{Np}{Neptunium}};
  \node[name=Pu, right of=Np, LanthanideActinide] {\SyntheticElementTextFormat{94}{244}{Pu}{Plutonium}};
  \node[name=Am, right of=Pu, LanthanideActinide] {\SyntheticElementTextFormat{95}{243}{Am}{Americium}};
  \node[name=Cm, right of=Am, LanthanideActinide] {\SyntheticElementTextFormat{96}{247}{Cm}{Curium}};
  \node[name=Bk, right of=Cm, LanthanideActinide] {\SyntheticElementTextFormat{97}{247}{Bk}{Berkelium}};
  \node[name=Cf, right of=Bk, LanthanideActinide] {\SyntheticElementTextFormat{98}{251}{Cf}{Californium}};
  \node[name=Es, right of=Cf, LanthanideActinide] {\SyntheticElementTextFormat{99}{252}{Es}{Einsteinium}};
  \node[name=Fm, right of=Es, LanthanideActinide] {\SyntheticElementTextFormat{100}{257}{Fm}{Fermium}};
  \node[name=Md, right of=Fm, LanthanideActinide] {\SyntheticElementTextFormat{101}{258}{Md}{Mendelevium}};
  \node[name=No, right of=Md, LanthanideActinide] {\SyntheticElementTextFormat{102}{259}{No}{Nobelium}};
  \node[name=Lr, right of=No, LanthanideActinide] {\SyntheticElementTextFormat{103}{262}{Lr}{Lawrencium}};

%%% Draw dotted lines connecting Lanthanide breakout to main table
%  \draw (LaLu.north west) edge[dotted] (La.north west)
%        (LaLu.north east) edge[dotted] (Lu.north east)
%        (LaLu.south west) edge[dotted] (La.south west)
%        (LaLu.south east) edge[dotted] (Lu.south east);
%%% Draw dotted lines connecting Actinide breakout to main table
%  \draw (AcLr.north west) edge[dotted] (Ac.north west)
%        (AcLr.north east) edge[dotted] (Lr.north east)
%        (AcLr.south west) edge[dotted] (Ac.south west)
%        (AcLr.south east) edge[dotted] (Lr.south east);

%% Legend

%  \node at ($(La.north -| Fr.west) + (5em,-15em)$) [name=elementLegend, Element, fill=white]
%    {\NaturalElementTextFormat{Z}{mass}{Symbol}{Name}};
      \node[yshift=10cm] at ($(Mn.north -| Fe.west) + (5em,-15em)$) [name=elementLegend, Element, fill=white]
    {\NaturalElementTextFormat{Z}{mass}{Symbol}{Name}};
%  \node[Element, fill=white, right of=elementLegend, xshift=1em]
%    {\SyntheticElementTextFormat{}{}{man-made}{}} ;
\begin{scope}[xshift =6cm]
  \draw[black, AlkaliMetalFill] (0,0) rectangle +(1em, 1em) node[right, yshift=-1ex]{Alkali Metal};
  \draw[black, AlkalineEarthMetalFill] (0,-2em) rectangle +(1em, 1em) node[right, yshift=-1ex]{Alkaline Earth Metal};
  \draw[black, MetalFill] (0,-4em) rectangle +(1em, 1em) node[right, yshift=-1ex]{Transition metals};
  \draw[black, MetalloidFill] (0,-6em)  rectangle +(1em, 1em) node[right, yshift=-1ex]{Metalloid};
%  \draw[black, NonmetalFill] (0,-8em)  rectangle +(1em, 1em) node[right, yshift=-1ex]{Non-metal};
  \draw[black, ChalcogenFill] (0,-8em)  rectangle +(1em, 1em) node[right, yshift=-1ex]{Chalcogen};
  \draw[black, HalogenFill] (0,-10em)  rectangle +(1em, 1em) node[right, yshift=-1ex]{Halogen};
  \draw[black, NobleGasFill] (0,-12em)  rectangle +(1em, 1em) node[right, yshift=-1ex]{Noble Gas};
  \draw[black, LanthanideActinideFill] (0,-14em)   rectangle +(1em, 1em) node[right, yshift=-1ex]{Lanthanide/Actinide};
\end{scope} 


%%% Diagram Title
%  \node at (H.west -| Fe.north) [name=diagramTitle, TitleLabel]
%    {Periodic Table of Chemical Elements};

%  \draw[black, AlkaliMetalFill] ($(La.north -| Fr.west) + (1em,-0.0em)$)
%    rectangle +(1em, 1em) node[right, yshift=-1ex]{Alkali Metal};
%  \draw[black, AlkalineEarthMetalFill] ($(La.north -| Fr.west) + (1em,-1.5em)$)
%    rectangle +(1em, 1em) node[right, yshift=-1ex]{Alkaline Earth Metal};
%  \draw[black, MetalFill] ($(La.north -| Fr.west) + (1em,-3.0em)$)
%    rectangle +(1em, 1em) node[right, yshift=-1ex]{Metal};
%  \draw[black, MetalloidFill] ($(La.north -| Fr.west) + (1em,-4.5em)$)
%    rectangle +(1em, 1em) node[right, yshift=-1ex]{Metalloid};
%  \draw[black, NonmetalFill] ($(La.north -| Fr.west) + (1em,-6.0em)$)
%    rectangle +(1em, 1em) node[right, yshift=-1ex]{Non-metal};
%  \draw[black, HalogenFill] ($(La.north -| Fr.west) + (1em,-7.5em)$)
%    rectangle +(1em, 1em) node[right, yshift=-1ex]{Halogen};
%  \draw[black, NobleGasFill] ($(La.north -| Fr.west) + (1em,-9.0em)$)
%    rectangle +(1em, 1em) node[right, yshift=-1ex]{Noble Gas};
%  \draw[black, LanthanideActinideFill] ($(La.north -| Fr.west) + (1em,-10.5em)$)
%    rectangle +(1em, 1em) node[right, yshift=-1ex]{Lanthanide/Actinide};
%    
    
\end{tikzpicture}
   

}






\newcommand{\EmptyTable}{
\begin{tikzpicture}[font=\sffamily, scale=0.350, transform shape, rotate=0, baseline={(0,9cm)}]

%% Fill Color Styles
  \tikzstyle{ElementFill} = [fill=yellow!15]
  \tikzstyle{AlkaliMetalFill} = [fill=blue!55]
  \tikzstyle{AlkalineEarthMetalFill} = [fill=blue!40]
  \tikzstyle{MetalFill} = [fill=blue!25]
  \tikzstyle{MetalloidFill} = [fill=orange!25]
  \tikzstyle{NonmetalFill} = [fill=green!25]
  \tikzstyle{HalogenFill} = [fill=green!40]
  \tikzstyle{NobleGasFill} = [fill=green!55]
  \tikzstyle{NoFill} = [fill=none]
  \tikzstyle{LanthanideActinideFill} = [fill=purple!25]

%% Element Styles
  \tikzstyle{Element} = [draw=black, ElementFill,
    minimum width=2.75cm, minimum height=2.75cm, node distance=2.75cm]
  \tikzstyle{AlkaliMetal} = [Element, AlkaliMetalFill]
  \tikzstyle{AlkalineEarthMetal} = [Element, AlkalineEarthMetalFill]
  \tikzstyle{Metal} = [Element, MetalFill]
  \tikzstyle{Metalloid} = [Element, MetalloidFill]
  \tikzstyle{Nonmetal} = [Element, NonmetalFill]
  \tikzstyle{Halogen} = [Element, HalogenFill]
  \tikzstyle{NobleGas} = [Element, NobleGasFill]
  \tikzstyle{LanthanideActinide} = [Element, LanthanideActinideFill]
  \tikzstyle{PeriodLabel} = [font={\sffamily\LARGE}, node distance=2.0cm]
  \tikzstyle{GroupLabel} = [font={\sffamily\LARGE}, minimum width=2.75cm, node distance=2.0cm]
  \tikzstyle{TitleLabel} = [font={\sffamily\Huge\bfseries}]
  \tikzstyle{plain} = [Element, NoFill]

%% Group 1 - IA
  \node[name=H, plain] {};
  \node[name=Li, below of=H, plain] {};
    \node[name=Na, below of=Li, plain] {};
  \node[name=K, below of=Na, plain] {};
  \node[name=Rb, below of=K, plain] {};
  \node[name=Cs, below of=Rb, plain] {};
  \node[name=Fr, below of=Cs, plain] {};

%% Group 2 - IIA
  \node[name=Be, right of=Li, plain] {};
  \node[name=Mg, below of=Be, plain] {};
  \node[name=Ca, below of=Mg, plain] {};
  \node[name=Sr, below of=Ca, plain] {};
  \node[name=Ba, below of=Sr, plain] {};
  \node[name=Ra, below of=Ba, plain] {};

%% Group 3 - IIIB
  \node[name=Sc, right of=Ca, plain] {};
  \node[name=Y, below of=Sc, plain] {};
  \node[name=LaLu, below of=Y, plain] {};
  \node[name=AcLr, below of=LaLu, plain] {};

%% Group 4 - IVB
  \node[name=Ti, right of=Sc, plain] {};
  \node[name=Zr, below of=Ti, plain] {};
  \node[name=Hf, below of=Zr, plain] {};
  \node[name=Rf, below of=Hf, plain] {};

%% Group 5 - VB
  \node[name=V, right of=Ti, plain] {};
  \node[name=Nb, below of=V, plain] {};
  \node[name=Ta, below of=Nb, plain] {};
  \node[name=Db, below of=Ta, plain] {};

%% Group 6 - VIB
  \node[name=Cr, right of=V, plain] {};
  \node[name=Mo, below of=Cr, plain] {};
  \node[name=W, below of=Mo, plain] {};
  \node[name=Sg, below of=W, plain] {};

%% Group 7 - VIIB
  \node[name=Mn, right of=Cr, plain] {};
  \node[name=Tc, below of=Mn, plain] {};
  \node[name=Re, below of=Tc, plain] {};
  \node[name=Bh, below of=Re, plain] {};

%% Group 8 - VIIIB
  \node[name=Fe, right of=Mn, plain] {};
  \node[name=Ru, below of=Fe, plain] {};
  \node[name=Os, below of=Ru, plain] {};
  \node[name=Hs, below of=Os, plain] {};

%% Group 9 - VIIIB
  \node[name=Co, right of=Fe, plain] {};
  \node[name=Rh, below of=Co, plain] {};
  \node[name=Ir, below of=Rh, plain] {};
  \node[name=Mt, below of=Ir, plain] {};

%% Group 10 - VIIIB
  \node[name=Ni, right of=Co, plain] {};
  \node[name=Pd, below of=Ni, plain] {};
  \node[name=Pt, below of=Pd, plain] {};
  \node[name=Ds, below of=Pt, plain] {};

%% Group 11 - IB
  \node[name=Cu, right of=Ni, plain] {};
  \node[name=Ag, below of=Cu, plain] {};
  \node[name=Au, below of=Ag, plain] {};
  \node[name=Rg, below of=Au, plain] {};

%% Group 12 - IIB
  \node[name=Zn, right of=Cu, plain] {};
  \node[name=Cd, below of=Zn, plain] {};
  \node[name=Hg, below of=Cd, plain] {};
  \node[name=Uub, below of=Hg, plain] {};

%% Group 13 - IIIA
  \node[name=Ga, right of=Zn, plain] {};
    \node[name=Al, above of=Ga, plain] {};
  \node[name=B, above of=Al, plain] {};
  \node[name=In, below of=Ga, plain] {};
  \node[name=Tl, below of=In, plain] {};
  \node[name=Uut, below of=Tl, plain] {};

%% Group 14 - IVA
  \node[name=C, right of=B, plain] {};
  \node[name=Si, below of=C, plain] {};
  \node[name=Ge, below of=Si, plain] {};
  \node[name=Sn, below of=Ge, plain] {};
  \node[name=Pb, below of=Sn, plain] {};
  \node[name=Uuq, below of=Pb, plain] {};

%% Group 15 - VA
  \node[name=N, right of=C, plain] {};
  \node[name=P, below of=N, plain] {};
  \node[name=As, below of=P, plain] {};
  \node[name=Sb, below of=As, plain] {};
  \node[name=Bi, below of=Sb, plain] {};
  \node[name=Uup, below of=Bi, plain] {};

%% Group 16 - VIA
  \node[name=O, right of=N, plain] {};
  \node[name=S, below of=O, plain] {};
  \node[name=Se, below of=S, plain] {};
  \node[name=Te, below of=Se, plain] {};
  \node[name=Po, below of=Te, plain] {};
  \node[name=Uuh, below of=Po, plain] {};

%% Group 17 - VIIA
  \node[name=F, right of=O, plain] {};
  \node[name=Cl, below of=F, plain] {};
  \node[name=Br, below of=Cl, plain] {};
  \node[name=I, below of=Br, plain] {};
  \node[name=At, below of=I, plain] {};
  \node[name=Uus, below of=At, plain] {}; 

%% Group 18 - VIIIA
  \node[name=Ne, right of=F, plain] {};
  \node[name=He, above of=Ne, plain] {};
  \node[name=Ar, below of=Ne, plain] {};
  \node[name=Kr, below of=Ar, plain] {};
  \node[name=Xe, below of=Kr, plain] {};
  \node[name=Rn, below of=Xe, plain] {};
  \node[name=Uuo, below of=Rn, plain] {}; 

%% Period
  \node[name=Period1, left of=H, PeriodLabel] {1};
  \node[name=Period2, left of=Li, PeriodLabel] {2};
  \node[name=Period3, left of=Na, PeriodLabel] {3}; 
  \node[name=Period4, left of=K, PeriodLabel] {4}; 
  \node[name=Period5, left of=Rb, PeriodLabel] {5};
  \node[name=Period6, left of=Cs, PeriodLabel] {6};
  \node[name=Period7, left of=Fr, PeriodLabel] {7};

%% Group
  \node[name=Group1, above of=H, GroupLabel] {1 \hfill IA};
  \node[name=Group2, above of=Be, GroupLabel] {2 \hfill IIA};
  \node[name=Group3, above of=Sc, GroupLabel] {3 \hfill IIIB};
  \node[name=Group4, above of=Ti, GroupLabel] {4 \hfill IVB};
  \node[name=Group5, above of=V, GroupLabel] {5 \hfill VB};
  \node[name=Group6, above of=Cr, GroupLabel] {6 \hfill VIB};
  \node[name=Group7, above of=Mn, GroupLabel] {7 \hfill VIIB};
  \node[name=Group8, above of=Fe, GroupLabel] {8 \hfill VIIIB};
  \node[name=Group9, above of=Co, GroupLabel] {9 \hfill VIIIB};
  \node[name=Group10, above of=Ni, GroupLabel] {10 \hfill VIIIB};
  \node[name=Group11, above of=Cu, GroupLabel] {11 \hfill IB};
  \node[name=Group12, above of=Zn, GroupLabel] {12 \hfill IIB};
  \node[name=Group13, above of=B, GroupLabel] {13 \hfill IIIA};
  \node[name=Group14, above of=C, GroupLabel] {14 \hfill IVA};
  \node[name=Group15, above of=N, GroupLabel] {15 \hfill VA};
  \node[name=Group16, above of=O, GroupLabel] {16 \hfill VIA};
  \node[name=Group17, above of=F, GroupLabel] {17 \hfill VIIA};
  \node[name=Group18, above of=He, GroupLabel] {18 \hfill VIIIA};

%% Lanthanide
  \node[name=La, below of=Rf, plain, yshift=-1cm] {};
  \node[name=Ce, right of=La, plain] {};
  \node[name=Pr, right of=Ce, plain] {};
  \node[name=Nd, right of=Pr, plain] {};
  \node[name=Pm, right of=Nd, plain] {};
  \node[name=Sm, right of=Pm, plain] {};
  \node[name=Eu, right of=Sm, plain] {};
  \node[name=Gd, right of=Eu, plain] {};
  \node[name=Tb, right of=Gd, plain] {};
  \node[name=Dy, right of=Tb, plain] {};
  \node[name=Ho, right of=Dy, plain] {};
  \node[name=Er, right of=Ho, plain] {};
  \node[name=Tm, right of=Er, plain] {};
  \node[name=Yb, right of=Tm, plain] {};
  \node[name=Lu, right of=Yb, plain] {};

%% Actinide
  \node[name=Ac, below of=La, plain, yshift=-1cm]{};
  \node[name=Th, right of=Ac, plain] {};
  \node[name=Pa, right of=Th, plain] {};
  \node[name=U, right of=Pa, plain] {};
  \node[name=Np, right of=U, plain] {};
  \node[name=Pu, right of=Np, plain] {};
  \node[name=Am, right of=Pu, plain] {};
  \node[name=Cm, right of=Am, plain] {};
  \node[name=Bk, right of=Cm, plain] {};
  \node[name=Cf, right of=Bk, plain] {};
  \node[name=Es, right of=Cf, plain] {};
  \node[name=Fm, right of=Es, plain] {};
  \node[name=Md, right of=Fm, plain] {};
  \node[name=No, right of=Md, plain] {};
  \node[name=Lr, right of=No, plain] {};
\end{tikzpicture} }
%%%%%%%%%%%%%%%%%%%
%%%%%%%%%%%%%%%%%%%
% ticks
\newcommand{\pro}[1]{\item[\ding{52}] {\color{green}#1}}
\newcommand{\con}[1]{\item[\ding{56}] {\color{red}#1}}

\usepackage{fontawesome}
\usepackage{enumitem}
\usepackage{cancel}
\newcommand\Cancel[1]{\cancel{\ce{#1}}}

\usepackage{savesym}
\usepackage{tikzorbital}
\usepackage{chemfig}\savesymbol{orbital}



\usepackage{chemmacros}
\usechemmodule{redox,reactions}
\chemsetup{modules=orbital}
\usepackage{enumitem}
\SetLabelAlign{myright}{\hss\llap{$#1$}}
\newlist{where}{description}{1}
\setlist[where]{labelwidth=2cm,labelsep=1em,
                        leftmargin=!,align=myright,font=\normalfont}

\usepackage{modiagram}
\newenvironment{myfont}{\fontfamily{augie}\selectfont}{\par}
\usepackage{mathrsfs}
\definecolor{notepadrule}{RGB}{217,244,244}


%%%%%%%%%FRUITS
\usepackage{xcolor-material}
\usetikzlibrary{fit}
\tikzset{%
   apple/.pic={
\fill [MaterialBrown] (-1/8,0) 
  arc (180:120:1 and 3/2) coordinate [pos=3/5] (@)-- ++(1/6,-1/7) 
  arc (120:180:5/4 and 3/2) -- cycle;
\fill [MaterialLightGreen500] (0,-9/10) 
  .. controls ++(180:1/8) and ++(  0:1/4) .. (-1/3,  -1)
  .. controls ++(180:1/3) and ++(270:1/2) .. (  -1,   0)
  .. controls ++( 90:1/3) and ++(180:1/3) .. (-1/2, 3/4)
  .. controls ++(  0:1/8) and ++(135:1/8) .. (   0, 4/7)
  .. controls ++( 45:1/8) and ++(180:1/8) .. ( 1/2, 3/4)
  .. controls ++(  0:1/3) and ++( 90:1/3) .. (   1,   0)
  .. controls ++(270:1/2) and ++(  0:1/3) .. ( 1/3,  -1)
  .. controls ++(180:1/4) and ++(  0:1/8) .. cycle;
\fill [MaterialLightGreen600] (0, 4/7)
  .. controls ++( 45:1/8) and ++(180:1/8) .. ( 1/2, 3/4)
  .. controls ++(  0:1/3) and ++( 90:1/3) .. (   1,   0)
  .. controls ++(270:1/2) and ++(  0:1/3) .. ( 1/3,  -1)
  .. controls ++(180:1/4) and ++(  0:1/8) .. (   0,-9/10);
\fill [MaterialGreen500, shift={(@)}, rotate=-30] 
  (0,0) arc (45:135:3/4 and 3/5) arc (225:315:3/4 and 3/5);
\fill [MaterialGreen700, shift={(@)}, rotate=-30] 
  (0,0) arc (315:225:3/4 and 3/5) -- cycle;
},
  orange/.pic={
\fill [MaterialOrange500] (0,0) circle [radius=1];
\fill [MaterialOrange600] (0,0) -- (45:1) arc (45:-135:1) -- cycle;
\fill [MaterialOrange700, shift={(0,3/4)}] coordinate (@)
ellipse [x radius=1/4, y radius=1/8];
\begin{scope}
\clip (0,0) circle [radius=1];
\fill [MaterialOrange700, shift=(@)] (90:1/4 and 1/8) 
  \foreach \i [evaluate={\j=mod(\i,2)+1/4;}]in {0,...,12}{
  -- (90+\i*30:\j*3/4 and \j*3/8) } -- cycle;
\end{scope}
\fill [MaterialBrown] (-1/16, 3/4) -- ++(0,1/4) arc (180:0:1/16 and 1/32)
   -- ++(0,-1/4) arc (360:180:1/16 and 1/32) -- cycle;
\fill [MaterialGreen500, shift=(@), rotate=-150] 
  (0,0) arc (45:135:1/2 and 4/5) arc (225:315:1/2 and 3/5);
\fill [MaterialGreen700, shift=(@), rotate=-150] 
  (0,0) arc (45:135:1/2 and 4/5) -- cycle;
},
lemon/.pic={
\tikzset{rotate=-45}
\fill [MaterialYellow500] (0, 0) ellipse [x radius=1/3, y radius=  1];
\fill [MaterialYellow500] (0, 0) ellipse [x radius=3/4, y radius=7/8];
\fill [MaterialYellow600] (270:1/3 and   1) arc (270:450:1/3 and   1);
\fill [MaterialYellow600] (270:3/4 and 7/8) arc (270:450:3/4 and 7/8);
\begin{scope}[shift=(90:9/10), rotate=-135]
\fill [MaterialGreen500] 
  (0,0) arc (45:135:1/2 and 4/5) arc (225:315:1/2 and 3/5);
\fill [MaterialGreen700] 
  (0,0) arc (45:135:1/2 and 4/5) -- cycle;
\tikzset{rotate=90, scale=3/4}
\fill [MaterialGreen500] 
  (0,0) arc (45:135:1/2 and 4/5) arc (225:315:1/2 and 3/5);
\fill [MaterialGreen700] 
  (0,0) arc (45:135:1/2 and 4/5) -- cycle;
\end{scope}
},
cherry/.pic={
\foreach \i in {1,2}{
\tikzset{shift={(-1+\i*3/4, -3/5+\i/5)},scale=1/2, rotate=15-\i*10}
\fill [MaterialRed700] (0,19/20) 
  .. controls ++(180:1/8) and ++(  0:1/4) .. (-1/3,  1)
  .. controls ++(180:1/3) and ++( 90:1/2) .. (  -1,  0)
  .. controls ++(270:1/2) and ++(180:1/2) .. (   0, -1)
  .. controls ++(  0:1/2) and ++(270:1/2) .. (   1,  0)
  .. controls ++( 90:1/2) and ++(  0:1/3) .. ( 1/3,  1)
  .. controls ++(180:1/4) and ++(  0:1/8) .. cycle;
\fill [MaterialRed800] (0, -1)
  .. controls ++(  0:1/2) and ++(270:1/2) .. (   1,  0)
  .. controls ++( 90:1/2) and ++(  0:1/3) .. ( 1/3,  1)
  .. controls ++(180:1/4) and ++(  0:1/8) .. (   0,19/20) -- cycle;
\fill [MaterialRed900] (0,3/4) coordinate (@\i) 
  ellipse [x radius=1/4, y radius=1/8];
}
\fill [MaterialBrown]
  (1/4,11/8) -- (3/8,11/8) coordinate (@)
  .. controls ++(270:1/2) and ++(135:1/3) .. (@1)
  .. controls ++(135:1/2) and ++(270:1/2) .. cycle;
\fill [MaterialBrown]
  (1/4,11/8) -- (3/8,11/8)
  .. controls ++(315:1/2) and ++(45:1/2) .. (@2)
  .. controls ++(60:1/2) and ++(315:1/2) .. cycle;
\fill [MaterialGreen500, shift={(@)}, rotate=20] 
  (0,0) arc (45:135:3/4 and 3/5) arc (225:315:3/4 and 3/5);
\fill [MaterialGreen700, shift={(@)}, rotate=20] 
  (0,0) arc (315:225:3/4 and 3/5) -- cycle;
}}
%%%%%%%%%FRUITS
%%%% FANCY COLOR TABLE
\usepackage{tcolorbox}
\usepackage{tabularx}
\usepackage{array}
\usepackage{colortbl}
\tcbuselibrary{skins}
\newcolumntype{Y}{>{\raggedright\arraybackslash}X}

\tcbset{tab1/.style={fonttitle=\bfseries\large,fontupper=\normalsize\sffamily,
colback=yellow!10!white,colframe=red!75!black,colbacktitle=Salmon!40!white,
coltitle=black,center title,freelance,frame code={
\foreach \n in {north east,north west,south east,south west}
{\path [fill=red!75!black] (interior.\n) circle (3mm); };},}}
\tcbset{tab2/.style={enhanced,fonttitle=\bfseries,fontupper=\normalsize\sffamily,
colback=yellow!10!white,colframe=red!50!black,colbacktitle=Salmon!40!white,
coltitle=black,center title}}
%%%% FANCY COLOR TABLE


 \newcolumntype{L}{>{\hsize=7.5\hsize}X}



%%%%% DNA
\usepackage{carbohydrates}
\definesubmol\nobond{[,0.2,,,draw=none]}
\definesubmol{adenine}{-[:72,,,,draw=none]\textcolor{blue}{N}(-[:-90]H)*5(-*6(-N=-N=(-NH_2)-=)--N=-)}
\definesubmol{adenine2}{-[:72,0.4,,,draw=none]\textcolor{blue}{N}*5(-*6(-N=-N=(-NH_2)-=)--N=-)}
\definesubmol{guanine}{-[:72,,,,draw=none]\textcolor{blue}{N}(-[:-90]H)*5(-*6(-N=(-[:-45]NH_2)-N(-[:45]H)-(=O)-=)--N=-)}
\definesubmol{guanine2}{-[:72,0.4,,,draw=none]\textcolor{blue}{N}*5(-*6(-N=(-[:-45]NH_2)-N(-[:45]H)-(=O)-=)--N=-)}

\definesubmol{uracil}{-[:90,,,,draw=none]\textcolor{blue}{N}(-[:-90]H)*6(-(=O)-N(-[:45]H)-(=O)-=-)}
\definesubmol{uracil2}{-[:90,0.4,,,draw=none]\textcolor{blue}{N}*6(-(=O)-N(-[:45]H)-(=O)-=-)}

\definesubmol{thymine}{-[:90,0.4,,,draw=none]\textcolor{blue}{N}(-[:-90]H)*6(-(=O)-N(-[:45]H)-(=O)-(-[:145]CH_3)=-)}
\definesubmol{thymine2}{-[:90,0.4,,,draw=none]\textcolor{blue}{N}*6(-(=O)-N(-[:45]H)-(=O)-(-[:145]CH_3)=-)}

\definesubmol{cytosine}{-[:90,,,,draw=none]\textcolor{blue}{N}(-[:-90]H)*6(-(=O)-N=(-NH_2)-=-)}
\definesubmol{cytosine2}{-[:90,0.4,,,draw=none]\textcolor{blue}{N}*6(-(=O)-N=(-NH_2)-=-)}

\definesubmol{ribose}{HO-[:90,0.5,2]?-[:60,1.0](-[:90,1.1]\textcolor{orange}{OH})-[:160,1.2]\textcolor{red}{O}-[:-160,1.2](-[:90,0.7](-[:180,0.9]H\textcolor{pink}{O}))-[:-60, 1.0]?(-[:-90,0.5]OH)}
\definesubmol{deoxyribose}{-[:90,0.5,2,,,draw=none]?-[:60,1.0](-[:90,1.1]OH)-[:160,1.2]\textcolor{red}{O}-[:-160,1.2](-[:90,0.7](-[:180,0.9]H\textcolor{pink}{O}))-[:-60, 1.0]?(-[:-90,0.5]OH)}
\definesubmol{phosphate}{-[6,.6]O-[6,1.5]P(-[4,.8]HO)(=[6,0.8]O)-[,1.5]O}
\newcommand*\derivesubmol[4]{%
    \saveexpandmode\saveexploremode\expandarg\exploregroups
    \csname @\ifcat\relax\noexpand#2first\else second\fi oftwo\endcsname
        {\expandafter\StrSubstitute\@car#2\@nil}
        {\expandafter\StrSubstitute\csname CF@@#2\endcsname}
    {\@empty#3}{\@empty#4}[\temp@]%
    \csname @\ifcat\relax\noexpand#1first\else second\fi oftwo\endcsname
        {\expandafter\let\@car#1\@nil}
        {\expandafter\let\csname CF@@#1\endcsname}\temp@
    \restoreexpandmode\restoreexploremode
}
\definesubmol{rt1}{-[2]rt1}
\definesubmol{rt2}{-[6,.6]rt2}
\definesubmol{rt3}{-[6,.6]rt3}
\definesubmol{ribose2}{%
    H\textcolor{pink}{O}-[:0,0.8]-[:-90,0.8]%
        (%
            -[:25,1.176]\textcolor{red}{O}%
            -[:-25,1.176]%
        )%
    -[:-45,0.8]%
        (%
          (-[:-90,0.8]OH)-[0,,,,line width=1pt,shorten <=-.5pt,shorten >=-.5pt]%
              (!{rt2})%
          -[:45,0.8]%
              (!{rt1})%
         )%
}
\definesubmol{deoxyribose2}{%
    H\textcolor{pink}{O}-[:0,0.8]-[:-90,0.8]%
        (%
            -[:25,1.176]\textcolor{red}{O}%
            -[:-25,1.176]%
        )%
    -[:-45,0.8]%
        (%
          (-[:-90,0.8]OH)-[0,,,,line width=1pt,shorten <=-.5pt,shorten >=-.5pt]%
              ()%
          -[:45,0.8]%
              (!{rt1})%
         )%
}
%\derivesubmol{deoxyribose2}{ribose2}{(!{rt2})}{   }
\definesubmol{deoxyribosephosphate}{%
    P(-[:-90,0.8]OH)(-[:90,0.8]OH)(=[:180,0.8]O)-\textcolor{pink}{O}-[:0,0.8]-[:-90,0.8]%
        (%
            -[:25,1.176]\textcolor{red}{O}%
            -[:-25,1.176]%
        )%
    -[:-45,0.8]%
        (%
          (-[:-90,0.8]OH)-[0,,,,,shorten <=-.5pt,shorten >=-.5pt]%
          -[:45,0.8]%
              (!{rt1})%
         )%
}
\definesubmol{ribosephosphate}{%
    P(-[:-90,0.8]OH)(-[:90,0.8]OH)(=[:180,0.8]O)-\textcolor{pink}{O}-[:0,0.8]-[:-90,0.8]%
        (%
            -[:25,1.176]\textcolor{red}{O}%
            -[:-25,1.176]%
        )%
    -[:-45,0.8]%
        (%
          (-[:-90,0.8]OH)-[0,,,,,shorten <=-.5pt,shorten >=-.5pt]%
          (-[:-90,0.8]OH)-[:45,0.8]%
              (!{rt1})%
         )%
}
%%%%% DNA



 \newcommand{\cmark}{\ding{51}}%
\newcommand{\xmark}{\ding{55}}%
\usepackage{dirtytalk}


%%%%%%COLORS
\definecolor{olive}{rgb}{0.3, 0.4, .1}
\definecolor{fore}{RGB}{249,242,215}
\definecolor{back}{RGB}{51,51,51}
\definecolor{title}{RGB}{255,0,90}
\definecolor{dgreen}{rgb}{0.,0.6,0.}
\definecolor{gold}{rgb}{1.,0.84,0.}
\definecolor{JungleGreen}{cmyk}{0.99,0,0.52,0}
\definecolor{BlueGreen}{cmyk}{0.85,0,0.33,0}
\definecolor{RawSienna}{cmyk}{0,0.72,1,0.45}
\definecolor{Magenta}{cmyk}{0,1,0,0}
%%%%%%COLORS

 
 
 
 
\newcommand{\notes}[4][\empty]{%
    \noindent #4\vspace{10pt}\\
    \foreach \n in {1,...,#2}{%
        \ifthenelse{\equal{#1}{\empty}}
            {\rule{#3}{0.5pt}\\}
            {\rule{#3}{0.5pt}\vspace{#1}\\}
        }
}





%%%%%%%%FRUITS CHATER FANCY HEADER
%\newcommand{\chapterfruit}{cherry} % apple orange lemon  cherry
%\newcommand{\chaptercolorheading}{red} % apple orange lemon  cherry
% \pagestyle{fancy}
%\newcommand{\X}{\phantom{X}} % Filler to define baseline of empty circles
%\fancyhf{}
%\DeclareRobustCommand{\bul}{\begin{tikzpicture}[baseline={(current bounding box.south)}]
%\fill [\chaptercolorheading!40,anchor=base,baseline] circle (1mm);
%\end{tikzpicture}}
%\renewcommand{\chaptermark}[1]{\markboth{Chapter \thechapter~%
%\bul~#1}{}}
%\renewcommand{\sectionmark}[1]{\markright{Experiment \thesection~%
%\bul~#1}{}}
%\renewcommand\headrulewidth{0pt}
%\fancyhead[LO]{\leftmark}
%\fancyhead[RE]{\rightmark}
%\fancyhead[RO]
%{\begin{tikzpicture}[yshift=5cm,baseline, every node/.style={minimum size=8mm, anchor=base}]
%\path pic[scale=0.5] () at (-1,0) {cherry}
%node at (0,0) [shape=circle, fill=\chaptercolorheading!20] (0,0) {\X} 
%node at (1,0) [shape=circle, fill=\chaptercolorheading!40] (0,0) {\X} 
%node at (2,0) [shape=circle, fill=\chaptercolorheading!60] (0,0) {\thepage}; 
%\end{tikzpicture}}
%
%\fancyhead[LE]
%{\begin{tikzpicture}[yshift=5cm,baseline, every node/.style={minimum size=8mm, anchor=base}]
%\path node at (0,0) [shape=circle, fill=\chaptercolorheading!60] (0,0) {\thepage} 
%node at (1,0) [shape=circle, fill=\chaptercolorheading!40] (0,0) {\X} 
%node at (2,0) [shape=circle, fill=\chaptercolorheading!20] (0,0) {\X}
%pic[scale=0.5] () at (3,0) {\chapterfruit} 
%;
%\end{tikzpicture}}
%%%%%%%%FRUITS CHATER FANCY HEADER




\begin{document} 
\section{  Atoms and elements  }\clearpage\mbox{}\clearpage \subfile{./Atomsandelements.tex}

\end{document}



